\documentclass[11pt]{article}

\usepackage[margin=1in]{geometry}
\usepackage{amsmath,amssymb,amsthm,mathtools}
\usepackage{listings}
\usepackage{microtype}
\usepackage[hidelinks]{hyperref}

\newtheorem{theorem}{Theorem}
\newtheorem{lemma}[theorem]{Lemma}
\newtheorem{definition}[theorem]{Definition}

\title{A Tutorial on CJR:\\
Syntax, Semantics, and Separation Logic}
\author{Mechanized development in \texttt{theories/}}
\date{}

\lstdefinestyle{coq}{
  basicstyle=\ttfamily\small,
  columns=fullflexible,
  keepspaces=true,
  showstringspaces=false,
  breaklines=true,
  frame=single,
  framerule=0.3pt,
  xleftmargin=0.4em,
  xrightmargin=0.4em,
  aboveskip=0.8em,
  belowskip=0.8em,
  literate={Σ}{{$\Sigma$}}1
           {↦ₛ}{{$\mapsto_s$}}2
           {↦ₒ}{{$\mapsto_o$}}2
           {→}{{$\rightarrow$}}1
           {∀}{{$\forall$}}1
           {∃}{{$\exists$}}1
           {∗}{{$*$}}1
           {⌜}{{$\ulcorner$}}1
           {⌝}{{$\urcorner$}}1
}
\lstset{style=coq}

\newcommand{\cjr}{CJR}
\newcommand{\sem}[1]{[\![#1]\!]}
\newcommand{\step}{\longrightarrow}
\newcommand{\steps}{\longrightarrow^{*}}
\newcommand{\Stuck}{\mathsf{stuck}}
\newcommand{\WF}{\mathsf{wf}}
\newcommand{\dom}{\operatorname{dom}}
\newcommand{\lookup}[2]{#1(#2)}
\newcommand{\None}{\mathsf{None}}
\newcommand{\Some}{\mathsf{Some}}
\newcommand{\Unit}{\mathsf{()}}
\newcommand{\True}{\mathsf{true}}
\newcommand{\False}{\mathsf{false}}
\newcommand{\NULL}{\mathsf{null}}
\newcommand{\pts}{\mathrel{\mapsto}}
\newcommand{\ptss}{\mathrel{\mapsto_{s}}}
\newcommand{\ptsr}{\mathrel{\mapsto_{r}}}
\newcommand{\ptsb}{\mathrel{\mapsto_{b}}}
\newcommand{\ptso}{\mathrel{\mapsto_{o}}}
\newcommand{\sep}{\mathbin{*}}
\newcommand{\wand}{\mathbin{-\!\!*}}
\newcommand{\later}{\mathop{\triangleright}}
\newcommand{\Wpre}{\mathsf{wp}}
\newcommand{\pure}[1]{\ulcorner #1 \urcorner}

\begin{document}
\maketitle

\begin{abstract}
We present \cjr{} (Cangjie Radical), the performance-parity core of a Cangjie-like language: unboxed structs, word-addressed \texttt{CPointer} blocks, and final classes, together with \texttt{func}, \texttt{if}, and \texttt{while}.
This note is a reading guide to the Coq/Iris mechanization.
We give every concept twice---once as the definition that is actually checked, and once as ordinary mathematics with the same content---and we follow each program-logic rule by an account of what it is for.
Finally, we walk through the examples that exercise the rules in order, and we conclude with a linked-list development, an array library, and an array list, each with the usual equations of that data structure.
\end{abstract}

\section{What this development is}

We take as our starting point a deliberately small core.
\cjr{} is sequential, monomorphic, and first-order; it has code pointers, so that a vtable can be written by hand.
It has no generics, no inheritance, no \texttt{enum}/\texttt{match}, and no concurrency.
Features of a larger source language that fall outside this cut are treated as sugar that must lower into the core at the same time and space cost.
What follows defines that core.
We do not define the lowering, a garbage collector, or byte-level layout.

Safety, in this development, is not a type system.
The grammar \texttt{ty} records the types one intends to use, but the operational semantics is untyped.
A program is safe when a weakest-precondition proof exists; adequacy then says that, from the empty machine, execution does not get stuck.

The Coq sources live under \texttt{theories/}.
The language is an Iris evaluation-context language~\cite{iris-ground-up,iris-essence}, so the frame rule and the bind rule are Iris's, rather than those of a separate Hoare calculus.
The overall shape follows HeapLang, the example language shipped with the Iris Coq development~\cite{iris-coq}.

\section{Syntax}

\subsection{Concrete syntax}

The programs in this note are written in Cangjie, as specified by the language documentation~\cite{cangjie-docs}.
\cjr{} is precisely the fragment of that language that the mechanization covers.
The core in \texttt{theories/lang.v} is the lowering target, not a second surface syntax.
A multi-argument \texttt{func} lowers to a single argument---a struct of the arguments---because core application is unary.
A class that is not marked \texttt{open} cannot be inherited; that is the final class of the core.
\texttt{CPointer} reads, writes, and offsets are the \texttt{unsafe} methods \texttt{read}, \texttt{write}, and \texttt{+} from the standard library, and allocation is the \texttt{foreign} function \texttt{malloc} from the Cangjie--C interoperability chapter.

\begin{lstlisting}
// types used below: Int64, Float64, Bool, Unit, struct, class, CPointer<T>

func add(a: Int64, b: Int64): Int64 {
    a + b
}

let frozen: Int64 = 1
var n: Int64 = 1
while (1 <= n) {
    n = n - 1
}

if (n <= 0) {
    0
} else {
    n
}

struct Cell {
    public var x: Int64
    public init(x: Int64) { this.x = x }
}

class Counter {          // not open, so it has no subclasses
    public var value: Int64
    public init(value: Int64) { this.value = value }
}

foreign func malloc(size: UIntNative): CPointer<Unit>
foreign func free(ptr: CPointer<Unit>): Unit
\end{lstlisting}

\texttt{let} is immutable, and in the core it is substitution.
\texttt{var} is mutable, and it is a stack slot.
The last expression of a block is its value, and \texttt{()} is unit.

\subsection{From Cangjie to the Coq core}

We summarise the correspondence in the table below.
Each row pairs one Cangjie form with the constructor in \texttt{theories/lang.v} to which it lowers.
Field names become positions: the first field of a \texttt{struct} or \texttt{class} is index \(0\).
A call \texttt{f(a, b)} is \texttt{App} of \texttt{f} to a \texttt{StructV} of the arguments, because \texttt{App} takes one argument.
Forms that mention a string name of a \texttt{var} are rewritten by \texttt{subst\_var} into the administrative stack operations listed in the last rows.
Those operations have no spelling in Cangjie; they appear only after a \texttt{var} has been given a stack slot.

\begin{center}
\small
\begin{tabular}{@{}p{0.46\textwidth}p{0.48\textwidth}@{}}
\hline
Cangjie & Coq constructor \\
\hline
\texttt{Int64}, \texttt{Float64}, \texttt{Bool}, \texttt{Unit}
  & \texttt{TyInt}, \texttt{TyFloat}, \texttt{TyBool}, \texttt{TyUnit} \\
\texttt{struct S \{ ... \}}
  & \texttt{TyStruct} of the field types \\
\texttt{CPointer<T>}
  & \texttt{TyPtr} \\
\texttt{class C \{ ... \}} with no \texttt{open}
  & \texttt{TyObj} of the field count \\
\texttt{(T) -> U}
  & \texttt{TyFun} \\
\hline
\texttt{1}, \texttt{1.0}, \texttt{true}, \texttt{()}
  & \texttt{LitInt}, \texttt{LitFloat}, \texttt{LitBool}, \texttt{LitUnit} \\
a \texttt{CPointer} value, including null from \texttt{CPointer<T>()}
  & \texttt{LitPtr}, \texttt{LitNull} \\
a class instance
  & \texttt{LitObj} \\
\texttt{func f(x) \{ e \}}
  & \texttt{Rec} (some name, some parameter, body); a finished closure is \texttt{RecV} \\
a struct value such as \texttt{Cell(0)}
  & \texttt{StructV} \\
\hline
a name
  & \texttt{Var} \\
\texttt{let x = e1} followed by \texttt{e2}
  & \texttt{Let} \\
\texttt{var x = e1} followed by \texttt{e2}
  & \texttt{VarBind} \\
\texttt{x = e} for a \texttt{var x}
  & \texttt{Assign} \\
\texttt{if (e0) \{ e1 \} else \{ e2 \}}
  & \texttt{If} \\
\texttt{while (e1) \{ e2 \}}
  & \texttt{While} \\
two expressions in sequence, or a block
  & \texttt{Seq} \\
\texttt{+}, \texttt{-}, \texttt{<=}, \texttt{==} on \texttt{Int64}, \texttt{\&\&}
  & \texttt{PlusOp}, \texttt{MinusOp}, \texttt{LeOp}, \texttt{EqOp}, \texttt{AndOp} \\
\texttt{f(a)}
  & \texttt{App} \\
\hline
\texttt{Cell(e)}
  & \texttt{Struct}, evaluated left to right \\
\texttt{p.x} on a struct
  & \texttt{StructLoad} at the index of \texttt{x} \\
\texttt{p.x = e} on a \texttt{var} struct
  & \texttt{StructStore} of that variable and index \\
\texttt{Counter(e)}
  & \texttt{New} \\
\texttt{o.value}
  & \texttt{FieldLoad} \\
\texttt{o.value = e}
  & \texttt{FieldStore} \\
\hline
\texttt{malloc} of one word, then \texttt{CPointer<Int64>(...)}
  & \texttt{Alloc} of \texttt{1}, result \texttt{LitPtr} \\
\texttt{free(p)}
  & \texttt{Free} \\
\texttt{unsafe \{ p.read() \}}
  & \texttt{Load} \\
\texttt{unsafe \{ p.write(v) \}}
  & \texttt{Store} \\
\texttt{unsafe \{ p + i \}}
  & \texttt{Offset} \\
\hline
no Cangjie spelling; introduced for a \texttt{var} slot
  & \texttt{Pop}, \texttt{StackLoad}, \texttt{StackAssign}, \texttt{StackFieldStore}, and the literal \texttt{LitStack} \\
\hline
\end{tabular}
\end{center}

The rest of this section defines those constructors and the machine on which they run.

\subsection{Locations, identifiers, and types}

A location is an integer.
We wrap stack slots, raw words, and objects in distinct constructors so that the four ghost maps of the program logic cannot be confused.

\begin{lstlisting}
Definition loc : Type := Z.
Inductive stack_id := StackId (z : Z).
Inductive raw_id := RawId (z : Z).
Inductive obj_id := ObjId (z : Z).

Inductive ty :=
  | TyInt | TyFloat | TyBool | TyUnit
  | TyStruct (ts : list ty)
  | TyPtr (t : ty)
  | TyObj (nfields : nat)
  | TyFun (t1 t2 : ty).
\end{lstlisting}

Informally, locations are \(\mathbb{Z}\), and the three identifier sorts are the tagged copies \(\mathsf{StackId}(\ell)\), \(\mathsf{RawId}(\ell)\), and \(\mathsf{ObjId}(\ell)\).
The type grammar is
\[
\begin{aligned}
\tau ::=\;&
  \mathsf{Int64}
  \mid \mathsf{Float64}
  \mid \mathsf{Bool}
  \mid \mathsf{Unit}
  \mid \mathsf{struct}\,\overline{\tau}
  \mid \mathsf{CPointer}\,\tau
  \mid \mathsf{class}_{n}
  \mid \tau \to \tau.
\end{aligned}
\]
Here \(\mathsf{class}_{n}\) means a final class with \(n\) fields.
Nothing in the semantics consults \(\tau\).
The type is documentation of the fragment; the theorems that follow do not assume it.

\subsection{Values and expressions}

A binder is an optional name: \(\mathsf{None}\) is an ignored parameter.
Literals cover integers, a float payload (for which no float operations are defined), booleans, unit, raw pointers, null, object references, and stack addresses.
Values are literals, recursive closures, and unboxed struct tuples.
Expressions comprise the core forms, together with administrative forms (\texttt{Pop}, \texttt{StackLoad}, \texttt{StackAssign}, \texttt{StackFieldStore}) that are introduced when a mutable variable is allocated.

\begin{lstlisting}
Definition binder := option string.

Inductive base_lit :=
  | LitInt (n : Z) | LitFloat (n : Z) | LitBool (b : bool)
  | LitUnit | LitPtr (l : loc) | LitNull
  | LitObj (o : loc) | LitStack (l : loc).

Inductive val :=
  | LitV (l : base_lit)
  | RecV (f x : binder) (e : expr)
  | StructV (vs : list val)
with expr :=
  | Val (v : val) | Var (x : string)
  | Rec (f x : binder) (e : expr)
  | App (e1 e2 : expr)
  | Let (x : binder) (e1 e2 : expr)
  | VarBind (x : string) (e1 e2 : expr)
  | Assign (x : string) (e : expr)
  | Pop (l : stack_id) (e : expr)
  | StackLoad (e : expr)
  | StackAssign (e1 e2 : expr)
  | StackFieldStore (e1 : expr) (i : nat) (e2 : expr)
  | If (e0 e1 e2 : expr)
  | While (e1 e2 : expr)
  | Seq (e1 e2 : expr)
  | BinOp (op : bin_op) (e1 e2 : expr)
  | Alloc (e : expr) | Free (e : expr)
  | Load (e : expr) | Store (e1 e2 : expr)
  | Offset (e1 e2 : expr)
  | Struct (vs : list val) (es : list expr)
  | StructLoad (e : expr) (i : nat)
  | StructStore (x : string) (i : nat) (e : expr)
  | New (vs : list val) (es : list expr)
  | FieldLoad (e : expr) (i : nat)
  | FieldStore (e1 : expr) (i : nat) (e2 : expr).
\end{lstlisting}

In mathematics we write values
\[
v ::=\;
  n \mid f \mid b \mid \Unit \mid \mathsf{ptr}(\ell) \mid \NULL
  \mid \mathsf{obj}(\ell) \mid \mathsf{stk}(\ell)
  \mid \mathsf{rec}\,f\,x.\,e
  \mid \langle v_0,\ldots,v_{k-1}\rangle
\]
and we use the expression forms by the same names as the constructors.
\(\mathsf{BinOp}\) ranges over \(+\), \(-\), \(\le\), \(=\), and boolean conjunction.
Integer operations act on \(\mathsf{LitInt}\); conjunction acts on booleans.
Anything else makes \(\mathsf{bin\_op\_eval}\) return \(\None\), and the machine gets stuck.

Structs and objects are built by a list zipper.
\(\mathsf{Struct}\,\overline{v}\,\overline{e}\) has already evaluated the values \(\overline{v}\) and still has the expressions \(\overline{e}\) to evaluate, left to right.
\(\mathsf{New}\) is the same zipper for a final class.
When the expression list is empty, the zipper is done.

\texttt{Val} and \texttt{to\_val} identify which expressions are finished:

\begin{lstlisting}
Definition of_val (v : val) : expr := Val v.
Definition to_val (e : expr) : option val :=
  match e with Val v => Some v | _ => None end.
\end{lstlisting}

So the values of the language, as Iris sees them, are exactly the expressions \(\mathsf{Val}\,v\).

\subsection{Two kinds of binding}

We begin from a distinction that is already present in the surface language and then explain how the core encodes it.
Immutable \texttt{let} and function parameters are substitution; mutable \texttt{var} is not.
A \texttt{var} is given a fresh stack slot, and every use or assignment of that name is rewritten into an operation on the slot.
A bare struct value is still copied: mutation in place happens only through a \texttt{var} slot or through a stored field.

Those surface facts come from the Cangjie 1.0 user manual~\cite{cangjie-docs}.
The stack-slot encoding below is this note's model of them, not a spelling in the language.
\begin{itemize}
\item
  \textbf{Immutable vs mutable bindings.}
  The chapter \emph{Program Structure}
  (\url{https://docs.cangjie-lang.cn/en/docs/1.0.0/user_manual/source_en/basic_programming_concepts/program_structure.html})
  introduces the mutability modifiers \texttt{let} and \texttt{var}:
  a \texttt{let} variable is assigned once (initialization), whereas a \texttt{var} variable may be reassigned.
  The chapter \emph{Function Definition}
  (\url{https://docs.cangjie-lang.cn/en/docs/1.0.0/user_manual/source_en/function/define_functions.html})
  states that function parameters are immutable and cannot be assigned to in the function body.
\item
  \textbf{Struct copying.}
  The same \emph{Program Structure} chapter distinguishes value types (including \texttt{struct}) from reference types (including \texttt{class}):
  assigning a value-type variable copies the instance; later mutation of one copy does not affect the other.
  \emph{Creating a struct Instance}
  (\url{https://docs.cangjie-lang.cn/en/docs/1.0.0/user_manual/source_en/struct/create_instance.html})
  repeats that assignment and parameter passing copy a \texttt{struct}, and that changing fields through an instance requires both the struct variable and the fields to be \texttt{var}.
  The related page \emph{mut Function}
  (\url{https://docs.cangjie-lang.cn/en/docs/1.0.0/user_manual/source_en/struct/mut.html})
  adds that \texttt{mut} methods need a mutable (\texttt{var}) receiver; they cannot be called on a \texttt{let} struct.
\item
  \textbf{Stored class fields.}
  The chapter \emph{Classes}
  (\url{https://docs.cangjie-lang.cn/en/docs/1.0.0/user_manual/source_en/class_and_interface/class.html})
  states that \texttt{class} is a reference type: assignment does not copy the object, so field updates through one variable are visible through others that share it.
  That is the ``stored field'' case in the sentence above.
\end{itemize}
What follows---fresh stack slots, \texttt{subst\_var}, \texttt{StackLoad}/\texttt{StackAssign}/\texttt{StackFieldStore}, \texttt{LitStack}, and \texttt{Pop}---is the Coq core's administrative rewriting, not Cangjie surface syntax.

\begin{lstlisting}
Definition subst_binder (b : binder) (v : val) (e : expr) : expr :=
  match b with None => e | Some x => subst x v e end.
\end{lstlisting}

The function \(\mathsf{subst}\,x\,v\,e\) replaces free occurrences of \(x\) by the value \(v\), and it does not descend under a binder that captures \(x\).
On a value, literals are unchanged and the body of a closure is substituted only when \(x\) is not the function name or the parameter.

The mutable variant \(\mathsf{subst\_var}\,x\,\ell\,e\) does a different job.
A read \(\mathsf{Var}\,x\) becomes \(\mathsf{StackLoad}(\mathsf{stk}(\ell))\).
An assignment \(\mathsf{Assign}\,x\,e\) becomes \(\mathsf{StackAssign}(\mathsf{stk}(\ell), e')\).
A field write \(\mathsf{StructStore}\,x\,i\,e\) becomes \(\mathsf{StackFieldStore}(\mathsf{stk}(\ell), i, e')\).
Other occurrences of \(x\) are rewritten recursively, again respecting capture.
The slot address is itself a value, so later stack operations do not look up a string environment.

For a concrete example, the fragment \(\mathsf{var}\;x = 1\;\mathsf{in}\;(x = x + 1)\)
allocates a fresh slot \(\ell\), stores \(1\), and rewrites the body, reaching
\[
\mathsf{Pop}\,(\mathsf{StackId}(\ell))\,
  (\mathsf{StackAssign}(\mathsf{stk}(\ell),\,
    \mathsf{BinOp}\,{+}(\mathsf{StackLoad}(\mathsf{stk}(\ell)),\,1))).
\]
Informally that is \(\mathsf{pop}_{\ell}\,(\mathsf{stk}(\ell) := \mathsf{stk}(\ell) + 1)\):
the read of \(x\) is a load from the slot, the assignment writes the slot, and \(\mathsf{Pop}\) deletes it when the body finishes.

What is striking about this split is that it is the whole model of ``locals do not need a points-to fact until they are mutable.''
An immutable binding disappears into the syntax.
A mutable binding is a cell, and the cell has an owner.

\section{Machine state}

\subsection{Three regions, one counter}

The machine is a single record.
Locals and unboxed structs that are stored in \texttt{var} slots live in the stack map.
\texttt{CPointer} words live in the raw map, and each allocated block records its length in a second map on the same addresses.
Final class fields live in a map from \((\mathsf{ObjId}(\ell), i)\) to a value.
A single counter \(\mathsf{next}\) hands out fresh addresses for all three regions.

\begin{lstlisting}
Record state := MkState {
  cjr_stack  : gmap stack_id val;
  cjr_raw    : gmap raw_id val;
  cjr_blocks : gmap raw_id nat;
  cjr_obj    : gmap (obj_id * nat) val;
  cjr_next   : Z
}.
Definition state_init : state := inhabitant. (* all maps empty, next = 0 *)
\end{lstlisting}

Write a state as
\[
\sigma =
(\sigma_{s},\, \sigma_{r},\, \sigma_{b},\, \sigma_{o},\, \mathsf{next}(\sigma)).
\]
The initial state \(\sigma_0\) is \((\emptyset,\emptyset,\emptyset,\emptyset,0)\).

Indexing a pointer is arithmetic on the address, one word per element:
\[
\mathsf{ptr}(\ell) + i \;=\; \mathsf{ptr}(\ell+i).
\]
There is no byte offset and no padding.
A block of length \(n\) at base \(\ell\) occupies the words \(\ell,\ell+1,\ldots,\ell+n-1\).
Blocks are freed explicitly.
Class objects are not freed in this development: there is no \texttt{delete} for them, and there is no reachability argument.

Null is a pointer value, not an address in \(\sigma_{r}\).
Loading or storing through \(\NULL\) has no rule, so it is stuck.
An out-of-bounds or use-after-free access is stuck for the same reason: the lookup premise of load, store, or free fails.

\subsection{Well-formed states}

Not every 5-tuple is a state that the allocator is allowed to produce.
Well-formedness says that every address already in use lies strictly below the counter, that every block's words are present, and that live blocks do not overlap.

\begin{lstlisting}
Definition stack_below m n :=
  forall l, StackId l in dom m -> (l < n)%Z.
(* raw_below, block_below, obj_below are the same idea. *)

Definition block_cover blocks raw next :=
  forall l n, blocks !! RawId l = Some n ->
    (l + Z.of_nat n <= next)%Z /\
    forall i, i < n -> is_Some (raw !! RawId (l + Z.of_nat i)).

Definition blocks_disjoint blocks :=
  forall l1 n1 l2 n2,
    blocks !! RawId l1 = Some n1 ->
    blocks !! RawId l2 = Some n2 ->
    l1 <> l2 ->
    (l1 + Z.of_nat n1 <= l2 \/ l2 + Z.of_nat n2 <= l1)%Z.

Record state_wf (s : state) : Prop := {
  wf_stack : stack_below (cjr_stack s) (cjr_next s);
  wf_raw   : raw_below (cjr_raw s) (cjr_next s);
  wf_block : block_below (cjr_blocks s) (cjr_next s);
  wf_obj   : obj_below (cjr_obj s) (cjr_next s);
  wf_cover : block_cover (cjr_blocks s) (cjr_raw s) (cjr_next s);
  wf_disj  : blocks_disjoint (cjr_blocks s)
}.
\end{lstlisting}

In mathematics, \(\WF(\sigma)\) is the conjunction of:
\begin{align*}
&\forall \ell.\; \mathsf{StackId}(\ell)\in\dom(\sigma_{s}) \Rightarrow \ell < \mathsf{next}(\sigma),\\
&\text{and the same for raw addresses, block bases, and object addresses,}\\
&\forall \ell,n.\;
  \lookup{\sigma_{b}}{\mathsf{RawId}(\ell)}=\Some(n)
  \Rightarrow
  \ell+n \le \mathsf{next}(\sigma)
  \;\wedge\;
  \forall i<n.\;
  \lookup{\sigma_{r}}{\mathsf{RawId}(\ell+i)}\neq\None,\\
&\text{distinct live blocks have disjoint integer intervals.}
\end{align*}

The freshness lemmas are immediate from the first four clauses.
If \(\WF(\sigma)\) and \(\ell=\mathsf{next}(\sigma)\), then \(\ell\) is absent from the stack, from the raw heap, from the block map, and, at every field index, from the object heap.
Allocation is therefore the insertion of a key that is not already present.

The initial state is well formed: \texttt{state\_init\_wf}.
Each successful head step that the program logic uses is proved to preserve \(\WF\): stack allocation, deletion, and update; raw store; one-word allocation and freeing; class allocation; and field update.

One-word allocation is worth spelling out, because it is the case the logic rule talks about.
From base \(\ell=\mathsf{next}(\sigma)\),
\[
\begin{aligned}
\sigma_{r}' &= \{\mathsf{RawId}(\ell) \mapsto 0\} \cup \sigma_{r},\\
\sigma_{b}' &= \sigma_{b}[\mathsf{RawId}(\ell)\mapsto 1],\\
\mathsf{next}' &= \ell+1.
\end{aligned}
\]
The new word is in range, the new block is covered, and disjointness holds because the new interval \([\ell,\ell+1)\) starts at the old counter, past every previous address.
Freeing that block deletes the length and the word.
Other blocks keep their cells precisely because their intervals were disjoint from \([\ell,\ell+1)\).

Class allocation at \(o=\mathsf{next}(\sigma)\) inserts the fields \((\mathsf{ObjId}(o), i) \mapsto v_i\) for a list \(v_0,\ldots,v_{k-1}\) and advances the counter by one.
Those keys are fresh because every old object address is below \(o\).

\section{Small-step semantics}

\subsection{Head reduction}

A head step is a relation
\[
(e,\sigma) \step (e',\sigma')
\]
with no observations and no forked threads.
Both are hard-wired: \texttt{observation} is the empty type, and every rule produces an empty list of threads.
The language is deterministic on the nose.

\paragraph{Pure computation.}
These rules neither read nor write the heaps.
\begin{align*}
\mathsf{rec}\,f\,x.\,e
  &\step
  \mathsf{rec}\,f\,x.\,e
  &&\text{(as a value)}\\
(\mathsf{rec}\,f\,x.\,e)\,v
  &\step
  e[x\mapsto v][f\mapsto \mathsf{rec}\,f\,x.\,e]\\
\mathsf{let}\,x = v \mathsf{ in } e
  &\step
  e[x\mapsto v]\\
\mathsf{if}\;\True\;\mathsf{then}\;e_1\;\mathsf{else}\;e_2
  &\step e_1\\
\mathsf{if}\;\False\;\mathsf{then}\;e_1\;\mathsf{else}\;e_2
  &\step e_2\\
\mathsf{while}\;e_1\;\mathsf{do}\;e_2
  &\step
  \mathsf{if}\;e_1\;\mathsf{then}\;(e_2;\,\mathsf{while}\;e_1\;\mathsf{do}\;e_2)\;\mathsf{else}\;\Unit\\
v;\,e &\step e\\
v_1 \oplus v_2 &\step v
  &&\text{when }\sem{\oplus}(v_1,v_2)=\Some(v)\\
\mathsf{ptr}(\ell)+i &\step \mathsf{ptr}(\ell+i)\\
\langle \overline{v},\, v::\overline{e}\rangle
  &\step
  \langle \overline{v}{+}{+}[v],\, \overline{e}\rangle\\
\langle \overline{v}\rangle
  &\step
  \langle v_0,\ldots\rangle
  &&\text{(the struct value)}\\
\langle \overline{v}\rangle.i &\step v
  &&\text{when }v\text{ is the }i\text{th component}.
\end{align*}
The \texttt{while} rule duplicates the loop.
It is a head step, not an evaluation context, so the copies are syntax rather than a shared mutable closure.
The condition and the body are evaluated later, by ordinary contexts.

\paragraph{Stack.}
Let \(\ell=\mathsf{next}(\sigma)\).
\begin{align*}
\mathsf{var}\;x = v\;\mathsf{in}\;e
  &\step
  \mathsf{pop}_{\ell}\,(e\{x\mapsto \ell\})
  &&\text{and }\sigma_{s}(\mathsf{StackId}(\ell)):=v,\;\mathsf{next}:=\ell+1\\
\mathsf{pop}_{\ell}\,v
  &\step
  v
  &&\text{and delete the slot}\\
\mathsf{stk}(\ell)
  &\step
  v
  &&\text{if the slot holds }v\\
\mathsf{stk}(\ell) := v
  &\step
  \Unit
  &&\text{if the slot exists; it then holds }v\\
\mathsf{stk}(\ell).i := v
  &\step
  \Unit
  &&\text{if the slot holds a struct and index }i\text{ is in range}.
\end{align*}
The last rule replaces component \(i\) with \(\mathsf{set\_nth}\) and writes the struct back into the same slot.
No raw or object address is allocated.
The most important consequence is that a stack update can be framed against a raw points-to fact: the two maps are different components of \(\sigma\).

\paragraph{Raw heap.}
For \(n>0\), allocation of \(n\) words at \(\ell=\mathsf{next}(\sigma)\) returns \(\mathsf{ptr}(\ell)\), records length \(Z.\mathsf{to\_nat}(n)\), zeroes the words, and sets the counter to \(\ell+n\).
The proved program-logic rule specialises this to \(n=1\).
\begin{align*}
\mathsf{free}(\mathsf{ptr}(\ell))
  &\step \Unit
  &&\text{if a block length is recorded at }\ell;
     \text{ that block's words are deleted}\\
{!}\,\mathsf{ptr}(\ell)
  &\step v
  &&\text{if the word holds }v\\
\mathsf{ptr}(\ell) := v
  &\step \Unit
  &&\text{if the word exists; it then holds }v.
\end{align*}
Free deletes whatever length is stored.
The one-word reasoning rule asks separately for the length fact \(1\), so the client of that rule cannot free a longer block by mistake.

\paragraph{Objects.}
The object zipper steps like the struct zipper.
When it finishes, at \(o=\mathsf{next}(\sigma)\),
\[
\mathsf{new}(v_0,\ldots,v_{k-1})
\step
\mathsf{obj}(o)
\]
and the fields \((o,i)\mapsto v_i\) are added.
Field load and field store are lookup and update at \((\mathsf{ObjId}(o), i)\).
Store requires the field to exist already.
A missing field, like a missing word, is stuck.

\subsection{Evaluation contexts}

Head steps fire only at the root.
The context items say where a subexpression may be reduced.
Each item has one hole, and the hole is the expression currently being evaluated.
Arguments are evaluated left to right.
\texttt{if} evaluates only the condition.
\texttt{while} does not evaluate under itself: the unfolding rule is the head step above.

Representative items, written as contexts:
\[
\begin{aligned}
&E[\mathsf{let}\,x = [\,]\;\mathsf{in}\;e],\;
 E[\mathsf{var}\,x = [\,]\;\mathsf{in}\;e],\;
 E[\mathsf{pop}_{\ell}\,[\,]],\\
&E[[\,]\,e],\; E[v\,[\,]],\;
 E[\mathsf{if}\,[\,]\;\mathsf{then}\;e_1\;\mathsf{else}\;e_2],\;
 E[[\,];e],\\
&E[[\,]\oplus e],\; E[v\oplus [\,]],\;
 E[{!}\,[\,]],\; E[\mathsf{free}\,[\,]],\;
 E[[\,]:=e],\; E[v:=[\,]].
\end{aligned}
\]
Struct and \texttt{new} items evaluate the next expression in the zipper.
The Coq constructor list is \texttt{ectx\_item}, and \texttt{fill\_item} implements the grammar.

Iris then builds the language in the standard way.

\begin{lstlisting}
Lemma cjr_ectxi_mixin : EctxiLanguageMixin of_val to_val fill_item base_step.
Canonical Structure cjr_ectxi_lang := EctxiLanguage cjr_ectxi_mixin.
Canonical Structure cjr_ectx_lang := EctxLanguageOfEctxi cjr_ectxi_lang.
Canonical Structure cjr_lang := LanguageOfEctx cjr_ectx_lang.
\end{lstlisting}

A configuration steps by finding a context \(K\) and a head redex:
\[
K[e] \step K[e']
\quad\text{when}\quad
(e,\sigma)\step(e',\sigma').
\]
Filling a non-value never produces a value, a value never head-reduces, and a head reduction of a filled term implies that the hole was already a value.
Those are the mixin obligations; they are what let Iris treat \(K\) as a program context in the bind rule.

\subsection{Determinism}

\begin{lstlisting}
Lemma base_step_det e s k e' s' t k2 e2 s2 t2 :
  base_step e s k e' s' t ->
  base_step e s k2 e2 s2 t2 ->
  k = k2 /\ e' = e2 /\ s' = s2 /\ t = t2.
\end{lstlisting}

\begin{theorem}[Determinism of head reduction]
If \((e,\sigma)\step(e_1,\sigma_1)\) and \((e,\sigma)\step(e_2,\sigma_2)\), then \(e_1=e_2\) and \(\sigma_1=\sigma_2\).
Observations and forked threads, which are always empty, agree as well.
\end{theorem}

The proof is inversion of both derivations.
Where a rule reads a map or calls \(\sem{\oplus}\), the two lookups are the same, so the two results are the same.
There is no scheduler and no choice of fresh address other than \(\mathsf{next}(\sigma)\).

\section{Separation logic}

\subsection{What is owned}

The program logic does not talk about the whole machine.
It talks about the cells a piece of code is allowed to use.
Ownership is a ghost map---one authoritative map per region---with a points-to fact for each cell.
Full ownership is the fraction \(1\).

\begin{lstlisting}
Definition stack_mapsto (l : stack_id) (v : val) : iProp S :=
  ghost_map_elem stack_name l (DfracOwn 1) v.
Definition raw_mapsto (l : raw_id) (v : val) : iProp S :=
  ghost_map_elem raw_name l (DfracOwn 1) v.
Definition block_mapsto (l : raw_id) (n : nat) : iProp S :=
  ghost_map_elem block_name l (DfracOwn 1) n.
Definition obj_mapsto (o : obj_id) (i : nat) (v : val) : iProp S :=
  ghost_map_elem obj_name (o, i) (DfracOwn 1) v.
\end{lstlisting}

The notations in the sources are \(\ell \ptss v\), \(\ell \ptsr v\), \(\ell \ptsb n\), and \(o \ptso[i] v\).
We use four maps, rather than a single heap, because a stack slot, a raw word, a block length, and a class field are different kinds of cell.
A stack proof cannot accidentally be applied to a pointer.

The state interpretation ties the ghost maps to the concrete machine and remembers well-formedness:

\begin{lstlisting}
Definition cjr_state_interp (s : state) : iProp S :=
  ghost_map_auth stack_name 1 (cjr_stack s) *
  ghost_map_auth raw_name 1 (cjr_raw s) *
  ghost_map_auth block_name 1 (cjr_blocks s) *
  ghost_map_auth obj_name 1 (cjr_obj s) *
  pure (state_wf s).
\end{lstlisting}

At every step, Iris owns \(\mathsf{StateInterp}(\sigma)\).
A points-to fact is a fragment of one of those authoritative maps.
Updating a cell is updating both the fragment and the matching entry of \(\sigma\), and the \(\WF\) conjunct is re-established by the lemma for that step.

\subsection{Weakest precondition, in this instantiation}

The Iris instance is the standard one for a deterministic language with no thread pool of interest:
\[
\begin{aligned}
\mathsf{stateInterp}(\sigma) &= \mathsf{cjr\_state\_interp}(\sigma),\\
\mathsf{forkPost}(\_) &= \top,\\
\mathsf{numLatersPerStep}(\_) &= 0.
\end{aligned}
\]
The weakest precondition \(\Wpre\,e\,\{\Phi\}\) is Iris's.
For this tutorial the reading is:

\begin{definition}[Specification]
\(\Wpre\,e\,\{\Phi\}\) means: for every well-formed extension of the current ghost ownership, executing \(e\) neither gets stuck nor forks, and if it terminates with a value \(v\) then \(\Phi(v)\) holds.
Partial correctness is enough: a diverging loop satisfies a weakest precondition whose postcondition is never reached.
\end{definition}

The later \(\later\) and the credit \(\pounds 1\) appearing in the rules are Iris's accounting for steps, in the mode \texttt{HasLc}.
When a rule says
\[
\ell \ptsr v \wand \later(\pounds 1 \wand \ell \ptsr v \wand \Phi(v))
\vdash
\Wpre\,({!}\,\ell)\,\{\Phi\},
\]
the intended content is appealingly simple: ownership of the cell is enough to load it, and one gets the same cell back, together with the value that was there.
The later and the credit are discharged by the step.
In the example proofs they are thrown away with an introduction \(\mathtt{iIntros\ "> \_"}\).
They do not change the memory reasoning.

\subsection{Structural rules}

These three rules are why the logic is a separation logic rather than a Hoare logic with an explicit frame side-condition.
Bind and the pure-step lemma are proved in \texttt{primitive\_laws.v} by appealing to Iris.
The frame rule is Iris's own \texttt{wp\_frame\_l}; this development does not re-prove it.

\begin{lstlisting}
Lemma wp_bind K e Phi :
  wp e (fun v => wp (fill K (Val v)) Phi) |- wp (fill K e) Phi.
\end{lstlisting}

\begin{lemma}[Bind]
If \(K\) is an evaluation context, then
\[
\Wpre\,e\,\{\,v.\; \Wpre\,K[v]\,\{\Phi\}\,\}
\;\vdash\;
\Wpre\,K[e]\,\{\Phi\}.
\]
\end{lemma}

\noindent\textit{Intuition.}
To prove a specification of a term with a hole, we prove a specification of the hole, and we allow the postcondition of the hole to talk about the rest of the program.
Left-to-right evaluation is just a choice of \(K\).
No memory fact has to be mentioned.
Facts that the hole does not use stay outside this rule and are framed.

\begin{lemma}[Frame]
\[
P \sep \Wpre\,e\,\{\Phi\}
\;\vdash\;
\Wpre\,e\,\{\,v.\; P \sep \Phi(v)\,\}.
\]
\end{lemma}

\noindent\textit{Intuition.}
If a proof of \(e\) never mentions \(P\), then \(e\) cannot depend on the cells inside \(P\), and those cells survive \(e\) unchanged.
This is the rule that makes ``the other pointer is untouched'' a consequence of not putting that pointer into the callee's precondition.
It is not proved by induction on \cjr{} expressions.

\begin{lemma}[Consequence]
\[
\Wpre\,e\,\{\Phi\}
\;\wedge\;
(\forall v.\;\Phi(v)\wand \Psi(v))
\;\vdash\;
\Wpre\,e\,\{\Psi\}.
\]
\end{lemma}

\noindent\textit{Intuition.}
A stronger postcondition can be weakened.
The examples use this when a lemma such as \texttt{bump\_wp} returns both a pure result and a framed cell, and the continuation only needs one of them.
In Coq it is \texttt{wp\_wand} or \texttt{wp\_mono}.

\subsection{Pure rules}

A pure head step is one that is defined in every state, does not change the state, and has exactly one successor.
The generic lemma is \texttt{wp\_pure\_step}.
Each concrete rule below is an instance.
Where the step produces a value immediately, the postcondition is under a later, because that is how the generic lemma returns the credit.
Where the step produces another expression, the rule simply asks for a weakest precondition of that expression.

\paragraph{Closures and application.}
\begin{lstlisting}
Lemma wp_rec f x e Phi :
  later (pound 1 -*- Phi (RecV f x e)) |- wp (Rec f x e) Phi.
Lemma wp_app f x e v Phi :
  wp (subst_binder x v (subst_binder f (RecV f x e) e)) Phi
    |- wp (App (Val (RecV f x e)) (Val v)) Phi.
\end{lstlisting}
\[
\begin{aligned}
&\later(\Phi(\mathsf{rec}\,f\,x.\,e)) \vdash \Wpre\,(\mathsf{rec}\,f\,x.\,e)\,\{\Phi\},\\
&\Wpre\,e[x\mapsto v][f\mapsto \mathsf{rec}\,f\,x.\,e]\,\{\Phi\}
  \vdash
  \Wpre\,((\mathsf{rec}\,f\,x.\,e)\,v)\,\{\Phi\}.
\end{aligned}
\]
\textit{Intuition.}
Evaluating a function literal yields the closure as a value.
Calling it substitutes the argument and, if the function is recursive, the closure itself.
No heap fact is required: the code is data, and the data is the value in hand.
That is what lets a vtable be an ordinary struct of closures.

\paragraph{Let, if, while, sequencing, arithmetic, offset.}
\begin{align*}
\Wpre\,e[x\mapsto v]\,\{\Phi\}
  &\vdash \Wpre\,(\mathsf{let}\,x=v\;\mathsf{in}\;e)\,\{\Phi\}\\
\Wpre\,e_1\,\{\Phi\}
  &\vdash \Wpre\,(\mathsf{if}\;\True\;\mathsf{then}\;e_1\;\mathsf{else}\;e_2)\,\{\Phi\}\\
\Wpre\,e_2\,\{\Phi\}
  &\vdash \Wpre\,(\mathsf{if}\;\False\;\mathsf{then}\;e_1\;\mathsf{else}\;e_2)\,\{\Phi\}\\
\Wpre\,(\mathsf{if}\;e_1\;\mathsf{then}\;(e_2;\,\mathsf{while}\,e_1\,\mathsf{do}\,e_2)\;\mathsf{else}\;\Unit)\,\{\Phi\}
  &\vdash \Wpre\,(\mathsf{while}\,e_1\,\mathsf{do}\,e_2)\,\{\Phi\}\\
\Wpre\,e\,\{\Phi\}
  &\vdash \Wpre\,(v;e)\,\{\Phi\}\\
\sem{\oplus}(v_1,v_2)=\Some(v)
  &\Rightarrow
  \later(\Phi(v)) \vdash \Wpre\,(v_1\oplus v_2)\,\{\Phi\}\\
\later(\Phi(\mathsf{ptr}(\ell+i)))
  &\vdash \Wpre\,(\mathsf{ptr}(\ell)+i)\,\{\Phi\}.
\end{align*}

\textit{Intuition for \texttt{while}.}
The loop rule does not mention an invariant; it unfolds the loop once into a conditional.
The invariant is whatever assertion one has in hand when returning to the same syntactic loop.
If the condition is false, the unfolding takes the else branch and returns unit.
If it is true, the body runs and the same rule applies again.
Termination is not proved; the unfolding is used only as many times as the concrete execution needs.

\textit{Intuition for offset.}
Adding an index to a pointer is arithmetic.
It does not check that the resulting address lies in a block.
The check happens when someone loads or stores: without a points-to fact at the computed address, the load rule does not apply, and the machine would be stuck.
Bounds are therefore a matter of which points-to facts the proof owns, not of a side condition on \(+\).

\paragraph{Struct values.}
\[
\begin{aligned}
&\later(\Phi(\langle \overline{v}\rangle)) \vdash \Wpre\,\langle\overline{v}\rangle\,\{\Phi\},\\
&\overline{v}(i)=v \Rightarrow
  \later(\Phi(v)) \vdash \Wpre\,(\langle\overline{v}\rangle.i)\,\{\Phi\}.
\end{aligned}
\]
\textit{Intuition.}
A finished zipper is the unboxed tuple.
Reading a component copies it out.
The tuple is not a heap cell, so there is nothing to frame and nothing to free.
If the index is out of range, there is no rule.

Building a struct or an object from a non-empty expression list is also a pure step (\texttt{StructStepS}, \texttt{NewStepS}).
The development uses \texttt{wp\_pure\_step} for that step in the vtable example rather than a named lemma.

\subsection{Stack rules}

\begin{lstlisting}
Lemma wp_var_bind x v e Phi :
  (forall l, StackId l |->s v -*-
     wp (subst_var x l e)
        (fun w => exists v', StackId l |->s v' * Phi w))
  |- wp (VarBind x (Val v) e) Phi.

Lemma wp_stack_load l v Phi :
  l |->s v -*- later (pound 1 -*- l |->s v -*- Phi v)
    |- wp (StackLoad (stk l)) Phi.

Lemma wp_stack_assign l v w Phi :
  l |->s v -*- later (pound 1 -*- l |->s w -*- Phi unit)
    |- wp (StackAssign (stk l) (Val w)) Phi.

Lemma wp_stack_field l vs i w v Phi :
  vs !! i = Some w ->
  l |->s StructV vs -*-
    later (pound 1 -*- l |->s StructV (set_nth i v vs) -*- Phi unit)
    |- wp (StackFieldStore (stk l) i (Val v)) Phi.
\end{lstlisting}

In mathematics, with \(\ell\) a stack address:
\begin{align*}
&\bigl(\forall \ell.\;
   \ell \ptss v \wand
   \Wpre\,e\{x\mapsto\ell\}\,\{\,w.\;\exists v'.\; \ell \ptss v' \sep \Phi(w)\,\}\bigr)
  \vdash
  \Wpre\,(\mathsf{var}\;x=v\;\mathsf{in}\;e)\,\{\Phi\},\\
&\ell \ptss v \wand \later(\ell \ptss v \wand \Phi(v))
  \vdash \Wpre\,(\mathsf{stk}(\ell))\,\{\Phi\},\\
&\ell \ptss v \wand \later(\ell \ptss w \wand \Phi(\Unit))
  \vdash \Wpre\,(\mathsf{stk}(\ell):=w)\,\{\Phi\},\\
&\overline{v}(i)=w \Rightarrow
  \ell \ptss \langle\overline{v}\rangle
  \wand \later\bigl(\ell \ptss \langle\overline{v}[i\mapsto u]\rangle \wand \Phi(\Unit)\bigr)
  \vdash
  \Wpre\,(\mathsf{stk}(\ell).i := u)\,\{\Phi\}.
\end{align*}

\textit{Intuition for \texttt{var}.}
The rule allocates a fresh slot, hands exclusive ownership of it to the proof, and asks for a proof of the body after the name has been compiled into stack operations.
The contents may be changed.
At the end the slot must be given back, at whatever value it then has.
The \(\mathsf{pop}\) administrative step deletes it, so the caller does not see a stack fact.
This is lexical scope: the owner of the cell is the body, and the cell does not escape.

The existential is essential.
If the postcondition demanded the original value, an assignment inside the body could not be verified, because assignment replaces the points-to fact.

\textit{Intuition for load, assign, and field update.}
Load reads and returns the same fact.
Assignment exchanges the fact for one about the new value.
Field update is assignment of a struct, with one component changed, and it is only available when that component already exists.
A missing component is stuck, and there is no points-to fact that would let the rule fire.
None of these rules allocates on the raw heap or the object heap.

\subsection{Raw-heap rules}

\begin{lstlisting}
Lemma wp_alloc_one Phi :
  (forall l, RawId l |->r LitV (LitInt 0) -*-
             RawId l |->b 1 -*- Phi (LitV (LitPtr l)))
    |- wp (Alloc (Val (LitV (LitInt 1)))) Phi.

Lemma wp_free_one l v Phi :
  RawId l |->r v -*- RawId l |->b 1 -*- later (Phi unit)
    |- wp (Free (Val (LitV (LitPtr l)))) Phi.

Lemma wp_load l v Phi :
  RawId l |->r v -*- later (RawId l |->r v -*- Phi v)
    |- wp (Load (Val (LitV (LitPtr l)))) Phi.

Lemma wp_store l v w Phi :
  RawId l |->r v -*- later (RawId l |->r w -*- Phi unit)
    |- wp (Store (Val (LitV (LitPtr l))) (Val w)) Phi.
\end{lstlisting}

\begin{align*}
&\bigl(\forall \ell.\;
   \ell \ptsr 0 \sep \ell \ptsb 1 \wand \Phi(\mathsf{ptr}(\ell))\bigr)
  \vdash
  \Wpre\,(\mathsf{alloc}(1))\,\{\Phi\},\\
&\ell \ptsr v \sep \ell \ptsb 1 \wand \later(\Phi(\Unit))
  \vdash
  \Wpre\,(\mathsf{free}(\ell))\,\{\Phi\},\\
&\ell \ptsr v \wand \later(\ell \ptsr v \wand \Phi(v))
  \vdash
  \Wpre\,({!}\,\ell)\,\{\Phi\},\\
&\ell \ptsr v \wand \later(\ell \ptsr w \wand \Phi(\Unit))
  \vdash
  \Wpre\,(\ell := w)\,\{\Phi\}.
\end{align*}

\textit{Intuition.}
Allocation creates a cell that was not owned before, initialised to zero, together with a separate fact that a block of length one starts there.
The two facts have different jobs.
The points-to fact lets one load and store the word.
The block fact lets one free the block.
Free consumes both, so the word cannot be used again: there is no fact left to feed to the load rule, and the concrete word has been deleted, so the machine would get stuck if one tried.

Load does not disturb the cell.
Store keeps ownership and changes the value.
Neither rule needs the block fact.
A client may therefore load and store while the right to free sits unused in the context, and the frame rule keeps it there.

The operational rule for allocation accepts any positive length.
The verified rule above is the one-word instance, which the small examples use.
The array library uses the general rule: a points-to fact per word, plus one block fact for the base and the length.

\subsection{Object rules}

\begin{lstlisting}
Lemma wp_new vs Phi :
  (forall o, ([* list] i |-> v in vs, ObjId o |->o[i] v) -*-
             Phi (LitV (LitObj o)))
    |- wp (New vs []) Phi.

Lemma wp_field_load o i v Phi :
  ObjId o |->o[i] v -*- later (ObjId o |->o[i] v -*- Phi v)
    |- wp (FieldLoad (Val (LitV (LitObj o))) i) Phi.

Lemma wp_field_store o i v w Phi :
  ObjId o |->o[i] v -*- later (ObjId o |->o[i] w -*- Phi unit)
    |- wp (FieldStore (Val (LitV (LitObj o))) i (Val w)) Phi.
\end{lstlisting}

\begin{align*}
&\bigl(\forall o.\;
   \mathop{\ast}\limits_{i<k} o \ptso[i] v_i
   \wand \Phi(\mathsf{obj}(o))\bigr)
  \vdash
  \Wpre\,(\mathsf{new}(v_0,\ldots,v_{k-1}))\,\{\Phi\},\\
&o \ptso[i] v \wand \later(o \ptso[i] v \wand \Phi(v))
  \vdash
  \Wpre\,(o.i)\,\{\Phi\},\\
&o \ptso[i] v \wand \later(o \ptso[i] w \wand \Phi(\Unit))
  \vdash
  \Wpre\,(o.i := w)\,\{\Phi\}.
\end{align*}

\textit{Intuition.}
\texttt{new} of a finished argument list allocates a fresh object and returns one points-to fact per field, with the initialising values.
There is no block fact and no free rule: the object stays allocated.
Field load and field store are the same idea as raw load and store, on a different map.
Owning \(o \ptso[i] v\) says nothing about \(o \ptso[j] w\) for \(j\neq i\), and nothing about any raw word.
That is why a raw cell can be framed across a field update, and why two fields of one object can be reasoned about separately.

\subsection{Exclusivity}

\begin{lstlisting}
Lemma stack_mapsto_exclusive l v1 v2 :
  l |->s v1 -*- l |->s v2 |- False.
Lemma raw_mapsto_exclusive l v1 v2 :
  l |->r v1 -*- l |->r v2 |- False.
\end{lstlisting}

\begin{lemma}[Exclusive ownership]
For a stack address or a raw address \(\ell\), and any values \(v_1,v_2\),
\[
\ell \ptss v_1 \sep \ell \ptss v_2 \vdash \bot,
\qquad
\ell \ptsr v_1 \sep \ell \ptsr v_2 \vdash \bot.
\]
\end{lemma}

\textit{Intuition.}
Full ownership is unique.
The ghost-map algebra rejects two fragments of fraction \(1\) at the same key, even if they agree on the value.
This is the specification of non-aliasing.
If a swap is specified with two points-to assumptions, applying it to one address twice is impossible: the precondition is false.
The machine might still have a single cell there; the logic refuses to pretend there are two.

The same algebra applies to block facts and to object fields.
The development records the stack and raw cases because those are the ones the examples need to point at.

\section{Main theorems}

\subsection{Head reduction is deterministic}

Theorem~1 above is \texttt{base\_step\_det}.
Restated: the head reduction relation is a partial function of \((e,\sigma)\).
Combined with the context machinery, a configuration has at most one successor.
Stuckness is therefore unambiguous.
Either the term is a value, or exactly one head redex is ready and the rule applies, or no rule applies and the program is stuck.

\subsection{Well-formedness is an invariant of the verified steps}

The state interpretation contains \(\pure{\WF(\sigma)}\).
Each memory rule proves that the state after the step is well formed if the state before was.
The statements are \texttt{wf\_stack\_alloc}, \texttt{wf\_stack\_delete}, \texttt{wf\_stack\_upd}, \texttt{wf\_raw\_store}, \texttt{wf\_alloc\_one}, \texttt{wf\_free\_one}, \texttt{wf\_new}, and \texttt{wf\_field\_store}.

\begin{lemma}[Allocator invariant]
\(\WF(\sigma_0)\).
If \(\WF(\sigma)\) and a verified memory rule steps to \(\sigma'\), then \(\WF(\sigma')\).
\end{lemma}

The one-word cases are the ones used by the logic.
For example, freeing a block of length \(1\) at \(\ell\) preserves coverage of every other block because disjointness puts their words outside \(\{\ell\}\).

\subsection{Adequacy}

\begin{lstlisting}
Lemma cjr_adequacy S e (phi : val -> Prop) :
  (forall `{!cjrGS S} `{!invGS_gen HasLc S},
     |- wp e (fun v => pure (phi v))) ->
  adequate NotStuck e state_init (fun v _ => phi v).
\end{lstlisting}

\begin{theorem}[Adequacy]
Let \(\varphi\) be a pure predicate on values.
Suppose that in every Iris instantiation of the \cjr{} ghost state one can prove
\[
\Wpre\,e\,\{\,v.\; \pure{\varphi(v)}\,\}.
\]
Then, executing \(e\) from \(\sigma_0\), every reachable configuration can take a step or is a value, and every final value \(v\) satisfies \(\varphi(v)\).
\end{theorem}

\noindent\textit{What the proof does.}
It allocates four empty ghost maps, packages them as the \cjr{} ghost state, and observes \(\WF(\sigma_0)\).
Iris's adequacy theorem for weakest preconditions then applies.
The hypothesis is a proof that does not depend on a particular ghost name, so it can be used at the names chosen here.

\noindent\textit{What it does not say.}
It does not say that every closed expression is safe.
An expression with a free variable, a null dereference, or a load of an address nobody allocated has no such proof, and the theorem is silent.
It also does not extract a functional-correctness result beyond the pure \(\varphi\) that was proved.
If the postcondition is only ``the result is unit'', adequacy says the program returns unit without getting stuck; the interesting memory facts were used and consumed inside the proof.

The closed examples \texttt{sum\_one}, \texttt{bump\_keep}, \texttt{cell}, \texttt{class\_bump}, and \texttt{vtable\_prog} are instantiated at this theorem.
The other examples are open: their hypotheses are points-to facts, so they are stated as weakest preconditions rather than as runs from \(\sigma_0\).

\section{The examples}

Each example uses a feature that the previous one does not.
They are the files in \texttt{theories/examples/}, in the order below.
The programs are written in the Cangjie fragment above.
The Coq definitions are the core \texttt{expr} terms those files actually prove.

\subsection{Arithmetic: \texttt{arith.v}}

\paragraph{Program.}
\lstinputlisting{generated/cj/arith.cj}

\paragraph{Theorem.}
\[
\Wpre\,\mathsf{sum\_one}\,\{\,v.\; \pure{v=1}\,\}.
\]
From \(\sigma_0\), the program returns \(1\) and does not get stuck (\texttt{sum\_one\_adequate}).

\paragraph{What is new.}
Nothing is allocated on either heap.
The proof uses \texttt{var}, stack load and assignment, \texttt{while}, \texttt{if}, sequencing, and integer \(+\) and \(-\).
The loop condition is true once: \(n\) is \(1\), then \(0\).
On the true branch, \(s\) becomes \(1\) and \(n\) becomes \(0\).
On the false branch the loop returns unit, and the trailing read returns \(s\).

The stack facts at the end are \(s \ptss 1\) and \(n \ptss 0\).
They are given back to \texttt{wp\_var\_bind} through the existential, and \texttt{pop} deletes them.
The pure postcondition does not mention the stack.
That is the scope rule from the previous section, exercised on a loop that really mutates.

\subsection{Stack structs: \texttt{struct\_upd.v}}

\paragraph{Program.}
\lstinputlisting{generated/cj/struct_upd.cj}

\paragraph{Theorems.}
\[
\begin{aligned}
&\ell \ptsr v \vdash
  \Wpre\,\mathsf{bump}\,\{\,w.\; \pure{w=1} \sep \ell \ptsr v\,\},\\
&\Wpre\,\mathsf{bump\_keep}\,\{\,v.\; \pure{v=0}\,\}.
\end{aligned}
\]
The closed run from \(\sigma_0\) returns \(0\) (\texttt{bump\_adequate}).

\paragraph{What is new.}
The struct lives in a stack slot, not on the raw heap.
The field write is \texttt{wp\_stack\_field}: the slot changes from \(\langle 0\rangle\) to \(\langle 1\rangle\), and the read copies the component out with \texttt{wp\_struct\_load}.
No raw address is allocated by \(\mathsf{bump}\) itself.

The first theorem is the frame.
Its hypothesis is an arbitrary raw cell, and the postcondition returns that same cell.
The proof of the stack update never mentions it, so it cannot have been changed.
The closed program makes the same point operationally: the fresh raw word is \(0\) before \(\mathsf{bump}\) and \(0\) after, which is what the final load returns.
The field update's result \(1\) is established on the way and then discarded by sequencing, because the observation we keep is the unchanged raw word.

\subsection{One raw cell: \texttt{cell.v}}

\paragraph{Program.}
\lstinputlisting{generated/cj/cell.cj}

\paragraph{Theorem.}
\[
\Wpre\,\mathsf{cell}\,\{\,v.\; \pure{v=\Unit}\,\}.
\]
From \(\sigma_0\), the program returns unit and does not get stuck.

\paragraph{What is new.}
This is the first use of \(\ptsr\) and \(\ptsb\).
Allocation produces \(\ell \ptsr 0\) and \(\ell \ptsb 1\).
The store replaces the points-to fact by \(\ell \ptsr 7\).
The load reads \(7\) and returns the same fact.
The program then discards that \(7\) by sequencing and frees the block.
Free consumes both facts and returns unit.

The pure specification does not say ``the loaded value was \(7\)''; that fact is internal to the proof.
What adequacy reports is that the whole expression, whose value is the value of \texttt{free}, returns unit without a stuck load, store, or double free.
After \texttt{free}, no points-to fact remains, so a further load could not be verified.

\subsection{Swap and aliasing: \texttt{swap.v}}

\paragraph{Program.}
\lstinputlisting{generated/cj/swap.cj}

\paragraph{Theorems.}
\[
\begin{aligned}
&\ell_1 \ptsr v_1 \sep \ell_2 \ptsr v_2
  \vdash
  \Wpre\,\mathsf{swap}(\ell_1,\ell_2)\,
    \{\,\_.\; \ell_1 \ptsr v_2 \sep \ell_2 \ptsr v_1\,\},\\
&\ell \ptsr v_1 \sep \ell \ptsr v_2 \vdash \bot.
\end{aligned}
\]

\paragraph{What is new.}
The postcondition needs both cells.
The proof loads \(\ell_2\) while still owning \(\ell_1\), stores into \(\ell_1\), then stores the saved value into \(\ell_2\).
Each rule mentions only the cell it touches; the other fact sits in the context and is framed.

The second lemma is \texttt{raw\_mapsto\_exclusive}, re-exported as \texttt{swap\_aliased}.
It is not a second proof of the same swap.
It says that the swap specification cannot be instantiated at one address twice: the hypothesis of two points-to facts at that address is absurd.
An overlapping specification would be a different theorem---for example a single cell updated twice---and it is not claimed here.

\subsection{Indexed blocks: \texttt{sum\_block.v}}

\paragraph{Program.}
\lstinputlisting{generated/cj/sum_block.cj}

\paragraph{Theorem.}
\[
p \ptsr 4 \sep p \ptsb 1
\vdash
\Wpre\,\mathsf{sum\_at}(p)\,
\bigl\{\,v.\;
  \pure{v=4}
  \sep p \ptsr 4
  \sep p \ptsb 1
\,\bigr\}.
\]

\paragraph{What is new.}
The address that is loaded is not the base written in the source; it is the value of \(\mathsf{Offset}\), that is \(p+i\).
On the only iteration, \(i\) is \(0\), and \(p+0=p\), so the load is justified by the points-to fact at \(p\).
The addition itself is the pure offset rule; it does not consume the points-to fact.
The loop then increments \(i\) to \(1\), the condition \(i\le 0\) fails, and the sum \(4\) is returned.

The block-length fact is never used by the load, and it is still there at the end.
That is the split between ``I may read this word'' and ``I may free this block''.
The example does not free.
A load of \(p+1\) would not be provable from these hypotheses: there is no points-to fact at \(p+1\), and a one-word block does not contain that word.

The loop invariant, in the ordinary sense, is the assertion carried across the unfolding:
\(i\) and \(s\) are stack cells, \(s\) holds the sum of the prefix already visited, and the raw word is unchanged.
For a one-iteration loop the prefix is either empty or the single element \(4\).

\subsection{Final classes: \texttt{class\_upd.v}}

\paragraph{Program.}
\lstinputlisting{generated/cj/class_upd.cj}

\paragraph{Theorem.}
\[
\Wpre\,\mathsf{class\_bump}\,\{\,v.\; \pure{v=0}\,\}.
\]
From \(\sigma_0\), the program returns \(0\) and does not get stuck.

\paragraph{What is new.}
\texttt{new} allocates on the object map, not the raw map.
The client receives \(o \ptso[0] 0\).
The field store changes it to \(o \ptso[0] 5\), and the field load reads \(5\).
The raw word allocated first is \(0\), is never stored, and is what the program returns.

So the closed observation is again the framed cell: object update does not change the raw heap.
The proof makes the separation concrete.
After \texttt{wp\_alloc\_one}, the raw points-to fact stays in the context through \texttt{wp\_new}, \texttt{wp\_field\_store}, and \texttt{wp\_field\_load}.
None of those lemmas takes a raw fact as a hypothesis.
The final \texttt{wp\_load} is what consumes it.

\subsection{Framing across a call: \texttt{frame\_call.v}}

\paragraph{Program.}
\lstinputlisting{generated/cj/frame_call.cj}

\paragraph{Theorem.}
\[
p \ptsr v \sep q \ptsr 3
\vdash
\Wpre\,\mathsf{write\_call}(p,q)\,
\bigl\{\,w.\;
  \pure{w=3}
  \sep p \ptsr 9
  \sep q \ptsr 3
\,\bigr\}.
\]

\paragraph{What is new.}
The callee is a closure.
\texttt{wp\_rec} makes the closure, \texttt{wp\_app} substitutes the pointer \(p\) for \(x\), and \texttt{wp\_store} needs only \(p \ptsr v\).
The fact \(q \ptsr 3\) is not in the callee's precondition.
It is framed across the call, then used by the load, which returns \(3\).

This is the frame rule applied at a function boundary.
Nothing about \texttt{App} is special: substitution puts a value in the body, and the body is an ordinary store.
The guarantee that \(q\) is unchanged is simply the guarantee that the store's proof did not receive \(q\).

\subsection{A hand-written vtable: \texttt{vtable.v}}

\paragraph{Program.}
\lstinputlisting{generated/cj/vtable.cj}

\paragraph{Theorem.}
\[
\Wpre\,\mathsf{vtable\_prog}\,\{\,v.\; \pure{v=1}\,\}.
\]
From \(\sigma_0\), the program returns \(1\) and does not get stuck.

\paragraph{What is new.}
The function is a value, the struct is a tuple holding that value, and the call is an ordinary application of a component that was loaded out of the tuple.
There is no method table in the machine and no dynamic dispatch rule.
A ``method'' is a function whose first argument is the object, and a ``vtable'' is a struct the programmer assembled.

The proof follows that reading.
First, the closure is formed by \texttt{wp\_rec}.
Second, the struct zipper evaluates the closure and becomes \(\langle \mathit{inc}\rangle\) by the pure struct rules.
\texttt{new} yields \(o \ptso[0] 0\).
\texttt{wp\_struct\_load} copies the closure out; this does not touch the object.
\texttt{wp\_app} substitutes the object reference into the body.
The body loads field \(0\), adds one, and stores \(1\).
Finally, the trailing field load reads that \(1\), which is the value of the whole program.

All three regions of the story meet: a stack-free unboxed struct of code, a code pointer, and an object field.
The raw heap is not used.
The earlier examples already showed that leaving it unused is safe.

\section{A linked list}

The file \texttt{theories/examples/list.v} is a self-contained implementation of a singly linked list, together with the equations that say the implementation really is a list.
It is larger than the earlier examples because the specification is not a single pure result: it is a predicate on the object heap, and the programs are functions that allocate and read nodes.

\subsection{Why a node is an object with a boolean tag}

A textbook list cell is a pair of a head and a pointer to the rest, and the empty list is the null pointer.
\cjr{} has null, and loading through null is stuck, which is the right machine behaviour.
What it does not have is an equality test on pointers or on object references.
\texttt{EqOp} compares integers only.
A program therefore cannot branch on ``is this pointer null?''.

The sticking point is resolved by using a non-\texttt{open} class with three fields, and a boolean the language can already branch on.

\begin{lstlisting}
Definition nil_body : expr :=
  Let (Some "n")
    (New [LitV (LitBool false); LitV (LitInt 0); LitV LitUnit] [])
    (Seq (FieldStore (Var "n") 2 (Var "n")) (Var "n")).

Definition cons_body : expr :=
  Let (Some "x") (StructLoad (Var "a") 0)
    (Let (Some "xs") (StructLoad (Var "a") 1)
      (New [] [Val (LitV (LitBool true)); Var "x"; Var "xs"])).
\end{lstlisting}

Field \(0\) is the tag: \(\False\) means empty, \(\True\) means cons.
Field \(1\) is the head integer, which is \(0\) and unused when the tag is false.
Field \(2\) is the tail.
On an empty node it is that same node; on a cons node it is the value that represents the rest of the list, itself either an empty node or another cons node.
The allocation writes unit into the field and then stores the new object there before \texttt{nil} returns.
The points-to fact owns the field.
The value in the field is a pointer to the object, so the fact does not own the object a second time, and the predicate does not unfold.
Objects are not freed, so a cons cell stays in the heap after its tail has been read.
That matches the language: there is a \texttt{new} rule and no rule that deletes an object.

In Cangjie the operations are ordinary functions.
Core application takes one argument, so \texttt{cons(x, xs)} lowers to a call on the struct of those two arguments.
The other functions take the node.

\lstinputlisting{generated/cj/list.cj}

A Cangjie constructor cannot receive the object it is building, so the printed class takes an optional tail and stores it in a raw field.
Unit, the placeholder written by the allocation, is printed as \texttt{None}.
The property \texttt{next} reads that raw field back as a \texttt{Node}: the stored node when it is present, and the node itself when it is empty.
After \texttt{nil} returns, the raw tail has been set to the node.
\texttt{empty} does not read \texttt{next}.
The functions above are the Cangjie surface; \texttt{cons} of two arguments lowers to the unary core call on the struct \(\langle x, xs\rangle\).

\texttt{empty} returns the negation of the tag.
An empty node is empty; a cons node is not.
The conditional is the only test.
It does not inspect the tail pointer.

\subsection{The representation predicate}

On paper, a mathematical list \(xs\) is related to a value \(p\) by \(\mathsf{isList}(xs, p)\).

\begin{definition}[List representation]
\(\mathsf{isList}([], p)\) holds when \(p = \mathsf{obj}(o)\) and
\[
o \ptso[0] \False
\sep
o \ptso[1] 0
\sep
o \ptso[2] \mathsf{obj}(o).
\]
\(\mathsf{isList}(x :: xs, p)\) holds when \(p = \mathsf{obj}(o)\) and there is a tail value \(q\) with
\[
o \ptso[0] \True
\sep
o \ptso[1] x
\sep
o \ptso[2] q
\sep
\mathsf{isList}(xs, q).
\]
\end{definition}

The Coq predicate is the same recursion, as a separation-logic assertion.

\begin{lstlisting}
Fixpoint is_list (xs : list Z) (p : val) : iProp S :=
  match xs with
  | [] => exists o, pure (p = LitV (LitObj o)) *
      ObjId o |->o[0] LitV (LitBool false) *
      ObjId o |->o[1] LitV (LitInt 0) *
      ObjId o |->o[2] LitV (LitObj o)
  | x :: xs' => exists o nxt, pure (p = LitV (LitObj o)) *
      ObjId o |->o[0] LitV (LitBool true) *
      ObjId o |->o[1] LitV (LitInt x) *
      ObjId o |->o[2] nxt *
      is_list xs' nxt
  end.
\end{lstlisting}

\textit{What the separating conjunction is doing.}
The three fields of one node are three separate facts, so a proof may load the tag and still own the head and the tail.
The recursive fact \(\mathsf{isList}(xs, q)\) owns every node of the tail and nothing of the cons cell.
Those two regions are disjoint.
That is why ``the tail is the original list'' can be stated by returning the original pointer together with \(\mathsf{isList}(xs, p)\), while the cons cell's own three fields sit beside it rather than inside it.

A value that is not an object reference does not satisfy \(\mathsf{isList}\).
The lemma \texttt{is\_list\_obj} records the useful direction: if the predicate holds, the value is \(\mathsf{obj}(o)\) for some \(o\).
The proofs rewrite with that equality before they substitute, so an abstract pointer is a literal object reference by the time it is placed in a struct.
Substitution then computes.

\subsection{What each operation establishes}

\paragraph{Empty list.}
\[
\Wpre\,\mathsf{nil}()\,\{\,v.\; \mathsf{isList}([], v)\,\}.
\]
\texttt{wp\_new} of the three initialising values returns three points-to facts, with unit in the tail.
\texttt{wp\_field\_store} then writes the new object into that field, which is the empty clause.
Nothing else is allocated.
This is \texttt{nil\_spec}.

\paragraph{Cons.}
\[
\mathsf{isList}(xs, p)
\vdash
\Wpre\,\mathsf{cons}\,\langle x, p\rangle
\,\Bigl\{\,v.\;
  \exists o.\;
  v = \mathsf{obj}(o)
  \sep
  o \ptso[0] \True
  \sep
  o \ptso[1] x
  \sep
  o \ptso[2] p
  \sep
  \mathsf{isList}(xs, p)
\,\Bigr\}.
\]
The argument struct is unboxed.
\texttt{wp\_struct\_load} copies \(x\) and then \(p\) out of it; those loads do not touch the heap.
The body then allocates one object.
The third field is the original pointer \(p\), not a copy of the tail's nodes.
The fact \(\mathsf{isList}(xs, p)\) is framed across allocation: \texttt{wp\_new} does not ask for it, so it cannot be changed.
The postcondition is the unfolding of \(\mathsf{isList}(x :: xs, v)\) with the tail pointer pinned to the pointer that was passed in.
Pinning it is what makes the later inverse law an equality of values, not only an equality of mathematical lists up to some existential.
This is \texttt{cons\_spec}.

\paragraph{Head.}
\[
\mathsf{isList}(x :: xs, p)
\vdash
\Wpre\,\mathsf{head}\,p
\,\{\,v.\; \pure{v = x} \sep \mathsf{isList}(x :: xs, p)\,\}.
\]
The proof opens the cons clause, loads field \(1\), and puts the same facts back together.
The list predicate after the call is the list predicate before the call.
Head is a pure observation.
This is \texttt{head\_spec}.

\paragraph{Tail.}
\[
\begin{aligned}
&\mathsf{isList}(x :: xs, p) \vdash {}\\
&\quad\Wpre\,\mathsf{tail}\,p
\,\Bigl\{\,v.\;
  \mathsf{isList}(xs, v)
  \sep
  \exists o.\;
  p = \mathsf{obj}(o)
  \sep
  o \ptso[0] \True
  \sep
  o \ptso[1] x
  \sep
  o \ptso[2] v
\,\Bigr\}.
\end{aligned}
\]
Loading field \(2\) returns the tail pointer and gives the points-to fact back.
The recursive clause of \(\mathsf{isList}\) already owned \(\mathsf{isList}(xs, v)\) separately from the cons cell, so the postcondition splits the list into the cell that was examined and the list that remains.
It does not duplicate nodes.
This is \texttt{tail\_spec}.

\paragraph{Emptiness.}
\[
\begin{aligned}
\mathsf{isList}([], p)
  &\vdash
  \Wpre\,\mathsf{empty}\,p
  \,\{\,v.\; \pure{v = \True} \sep \mathsf{isList}([], p)\,\},\\
\mathsf{isList}(x :: xs, p)
  &\vdash
  \Wpre\,\mathsf{empty}\,p
  \,\{\,v.\; \pure{v = \False} \sep \mathsf{isList}(x :: xs, p)\,\}.
\end{aligned}
\]
Both proofs load the tag and take one branch of \texttt{if}.
The other branch is dead in that specification: a proof of the empty law never has to show anything about the cons branch, and conversely.
The predicate is restored, so emptiness is an observation rather than a destructive read.
These are \texttt{empty\_nil} and \texttt{empty\_cons}.

\subsection{Adding an element and removing it}

The program that the inverse law talks about is the composition of the two functions, not a fused implementation.

\begin{lstlisting}
Definition cons_then_tail (x : Z) (p : val) : expr :=
  Let (Some "node") (call_cons x p)
    (App (Rec None (Some "n") tail_body) (Var "node")).
\end{lstlisting}

\begin{theorem}[Cons then tail]
\[
\mathsf{isList}(xs, p)
\vdash
\Wpre\,(\mathsf{tail}\,(\mathsf{cons}\,\langle x, p\rangle))
\,\{\,v.\; \pure{v = p} \sep \mathsf{isList}(xs, v)\,\}.
\]
\end{theorem}

The Coq statement is \texttt{cons\_then\_tail\_spec}.
Read it as an equation of lists.
The value returned by removing the element just added is the original pointer, and that pointer still represents the original mathematical list.

The proof follows the program in order.
\begin{enumerate}
\item \texttt{cons\_spec} allocates the node and frames \(\mathsf{isList}(xs, p)\).
Its postcondition says field \(2\) of the new object is exactly \(p\), and the original list predicate is still owned.
\item The \texttt{let} substitutes that object for \texttt{node}.
\item \texttt{tail} is the function above.
It loads field \(2\).
The load rule returns the value that the points-to fact describes, which is \(p\), and it returns the points-to fact as well.
\item The cons cell's three fields are no longer needed for the equation.
They are discarded.
Iris assertions are affine, so ownership of leftover cells may be dropped.
The object remains allocated in the machine; the specification simply stops mentioning it.
What remains is \(\pure{v = p}\) and \(\mathsf{isList}(xs, p)\).
\end{enumerate}

Two companion readings are already in the lemmas above, and they are the other usual equations.
Head after cons returns \(x\) and still owns \(\mathsf{isList}(x :: xs, \cdot)\).
Emptiness distinguishes \([]\) from \(x :: xs\).
Together they say that the tag, the head field, and the tail field agree with the mathematical constructors, and that reading them does not damage the list those constructors describe.

The next section proves in-place reversal of an arbitrary finite list using this representation.
It opens one cons cell at a time, rewires its tail, and transfers that cell into a separate accumulator.
Append remains outside this development.

\section{In-place list reversal}
\label{sec:list-reversal}

The file \texttt{theories/examples/list\_rev.v} proves the classic iterative
reversal of a singly linked list. This example changes the links of existing
cons cells. It does not allocate replacement cons cells, copy the integers, or
change the tag and head fields. It does allocate one fresh empty sentinel at
entry, because the two list predicates in the loop invariant must own separate
empty nodes. The old sentinel remains allocated when the loop finishes.
The functional theorem describes the reversed contents; the proof also explains
which field changes at each iteration and why that mutation is safe.

\subsection{The representation being reused}

There is one list representation in the development: \texttt{list.is\_list}
from the preceding section. Reversal imports it rather than defining a second
predicate. Its two clauses are
\[
\begin{aligned}
\mathsf{isList}([],p) \equiv {}&
 \exists o.\;\pure{p=\mathsf{obj}(o)}
 \sep o\ptso[0]\False \sep o\ptso[1]0 \sep o\ptso[2]\mathsf{obj}(o),\\
\mathsf{isList}(x::v,p) \equiv {}&
 \exists o,q.\;\pure{p=\mathsf{obj}(o)}
 \sep o\ptso[0]\True \sep o\ptso[1]x \sep o\ptso[2]q
 \sep\mathsf{isList}(v,q).
\end{aligned}
\]
The predicate owns all three fields of each node, including the terminal empty
node. Its separating conjunction makes the recursive tail separate from the
current cons cell. A finite list described this way cannot contain a cycle of
cons cells or repeated ownership of a cons cell. The sentinel's self-link is
allowed because the empty clause owns its three fields directly and does not
recurse through that link.

This ownership is the permission to perform the reversal. Another logical
client cannot simultaneously hold full ownership of the same fields. An
ordinary program variable can still contain an alias to a node: the algorithm
mutates that node, so observations through such an alias change too. The
specification transfers the input list's ownership into ownership of the
returned list; it does not promise that the original head still denotes the
original sequence.

\subsection{The client specification}

The main statement is \texttt{reverse\_spec}:
\[
\mathsf{isList}(xs,p)\vdash
 \Wpre\,\mathsf{reverse}(p)
 \,\{\,r.\;\mathsf{isList}(\mathsf{rev}(xs),r)\,\}.
\]
Here \(xs\) is a mathematical list of integers and \(p\) is a runtime value.
The list predicate connects them. The output predicate establishes both the
order of all values and a well-formed owned list. There is no positive-length
precondition; the empty and singleton cases are included.

The theorem uses Iris's ordinary weakest precondition, hence partial
correctness and safety. It states the postcondition whenever execution returns
and rules out getting stuck under its ownership precondition. The loop proof
uses induction on the unprocessed finite list, which explains why each
iteration consumes one remaining element. This chapter does not claim a
separate total-correctness theorem, exact instruction count, physical pointer
identity theorem, or verified equivalence between the CJR semantics and the
Cangjie compiler. Executable tests complement the checked CJR theorem.

By framing, a separate assertion \(R\) can accompany the input and output:
\(\mathsf{isList}(xs,p)\sep R\) is transformed into
\(\mathsf{isList}(\mathsf{rev}(xs),r)\sep R\). The proof needs none of the
resources in \(R\); a disjoint second list is a concrete example.

\subsection{The iteration definition: \texttt{rev\_step}}

\begin{minipage}{\linewidth}
\lstinputlisting[firstline=61,lastline=65]{theories/examples/list_rev.v}
\end{minipage}

The immutable \texttt{let nxt} reads field \(2\) before the program overwrites
it. Without this saved value, the algorithm would lose access to the remaining
input. The field store then changes \texttt{cur.next} to the accumulator.
The assignments make the current node the new accumulator head and advance
\texttt{cur} to the saved tail. Their order matters: assigning \texttt{cur}
first would lose the reference needed for \texttt{acc := cur}.

Only field \(2\) is changed. The current node retains its true tag and its
integer. After the store, that node and the old accumulator describe one new
list with contents \(x::u\). The saved tail describes the remaining suffix
\(v'\). The logical split follows the physical change in links.

\subsection{The loop definition: \texttt{rev\_loop}}

\begin{minipage}{\linewidth}
\lstinputlisting[firstline=67,lastline=68]{theories/examples/list_rev.v}
\end{minipage}

The condition loads the current node's tag. True means a cons node and permits
one iteration; false means the empty sentinel and exits. The program does not
use null as the stopping condition. The guard never reads field \(2\) of the
empty node, so its self-link does not cause an extra iteration.

\subsection{The function body: \texttt{rev\_body}}

\begin{minipage}{\linewidth}
\lstinputlisting[firstline=70,lastline=74]{theories/examples/list_rev.v}
\end{minipage}

\texttt{call\_nil} allocates and initializes the fresh empty accumulator.
\texttt{VarBind} allocates mutable stack bindings for \texttt{acc} and
\texttt{cur}; \texttt{n} is the immutable function argument. The final expression
reads the accumulator after the loop.

\subsection{The call wrapper: \texttt{call\_reverse}}

\begin{minipage}{\linewidth}
\lstinputlisting[firstline=76,lastline=77]{theories/examples/list_rev.v}
\end{minipage}

\texttt{call\_reverse} packages the body as a closure with one immutable
argument \texttt{n}, then applies it to the supplied runtime value. The
wrapper makes the client theorem talk about one complete function call rather
than an open expression with a free argument.

For an input with \(n\) cons cells, the source algorithm traverses those cells
once and uses a constant number of local references. The only new list object
is the empty sentinel. This is the intended linear-time, constant-extra-space
algorithm; no resource-cost theorem is asserted by \texttt{reverse\_spec}.

\subsection{Loading a local: \texttt{sld}}

\begin{minipage}{\linewidth}
\lstinputlisting[firstline=57,lastline=57]{theories/examples/list_rev.v}
\end{minipage}

A \texttt{VarBind} replaces reads of a mutable variable with loads from a fresh
stack slot. \texttt{sld z} abbreviates this concrete stack-load expression.
Its argument is the slot's location, not the object reference stored there.
The iteration proof uses the stack points-to fact to obtain that reference
and then separately uses object-field ownership to access the node.

\subsection{Assigning a local: \texttt{sas}}

\begin{minipage}{\linewidth}
\lstinputlisting[firstline=58,lastline=59]{theories/examples/list_rev.v}
\end{minipage}

\texttt{sas z e} evaluates the expression \(e\) and writes the resulting value
into stack slot \(z\). It changes the local reference rather than any object
field. This distinction accounts for the three different writes in one
iteration: one object-link write and two local-reference writes.

\subsection{The iteration after substitution: \texttt{rev\_step\_at}}

\begin{minipage}{\linewidth}
\lstinputlisting[firstline=79,lastline=82]{theories/examples/list_rev.v}
\end{minipage}

After allocating the accumulator and current bindings at locations
\(l_a,l_c\), this is the iteration that actually executes in the CJR semantics.
It reads the saved tail through \texttt{sld lc}, stores \texttt{sld la} in the
current object's link, and performs the two local assignments using
\texttt{sas}. It exposes the evaluation contexts used in the ownership proof.
It is a proof-facing form of \texttt{rev\_step}, not a second algorithm.

\subsection{The loop after substitution: \texttt{rev\_loop\_at}}

\begin{minipage}{\linewidth}
\lstinputlisting[firstline=84,lastline=85]{theories/examples/list_rev.v}
\end{minipage}

The guard loads field \(0\) of the object referenced by the current stack
slot, and the body is the substituted iteration. The accumulator location
appears only in the body. Cangjie is printed from \texttt{rev\_body}, whose
local variables retain their source names; these administrative stack forms
are used to connect that source expression to primitive proof rules.

\subsection{Iteration substitution: statement and proof}

\begin{minipage}{\linewidth}
\lstinputlisting[firstline=87,lastline=91]{theories/examples/list_rev.v}
\end{minipage}

This equality connects the named iteration with code operating on two stack
slots. The proof reduces \texttt{subst\_var} structurally through the
expression constructors. \texttt{f\_equal} reduces equality of matching
constructors to equality of their arguments; the remaining equalities are
reflexive. This is an ordinary Coq equality between syntax trees. It assumes
no ownership and executes no command.

\subsection{Loop substitution: statement and proof}

\begin{minipage}{\linewidth}
\lstinputlisting[firstline=93,lastline=98]{theories/examples/list_rev.v}
\end{minipage}

This proof unfolds the loop, reduces substitution in the condition, and uses
\texttt{rev\_step\_at\_subst} for its body. Thus it connects the entire named
loop to the loop verified over concrete stack slots. The main function proof
uses the reverse order of the substitutions, which is definitionally the same
for these two distinct variable names.

\subsection{The loop invariant: \texttt{rev\_inv}}

\begin{minipage}{\linewidth}
\lstinputlisting[firstline=100,lastline=106]{theories/examples/list_rev.v}
\end{minipage}

The witnesses \(u,v\) are respectively the contents of the accumulator and
the unprocessed suffix. The witnesses \(a,c\) are their runtime references.
The assertion says
\[
\exists u,v,a,c.\;
 \pure{\mathsf{rev}(u)\mathbin{+\mkern-6mu+}v=xs}
 \sep l_a\ptss a \sep l_c\ptss c
 \sep\mathsf{isList}(v,c)\sep\mathsf{isList}(u,a).
\]
The accumulator already has reversed order. Reversing it mathematically
recovers the processed prefix of the original input. Appending the remaining
suffix must recover the entire original list. The stack facts describe the
current values of the two mutable bindings, and the two list predicates give
separate ownership of both chains.

For \(xs=[1,2,3]\), the successive pairs \((u,v)\) are
\(([],[1,2,3])\), \(([1],[2,3])\), \(([2,1],[3])\), and
\(([3,2,1],[])\). For instance,
\(\mathsf{rev}([2,1])\mathbin{+\mkern-6mu+}[3]=[1,2,3]\).
At entry, \(u=[]\) and \(v=xs\). At exit, \(v=[]\), and the equation gives
\(u=\mathsf{rev}(xs)\).

Fresh allocation of the accumulator sentinel is essential to this particular
invariant. Reusing the original terminal node would demand its full ownership
inside both list predicates. The fresh sentinel lets the two complete chains
remain separate throughout the traversal. At exit the original sentinel is
owned by \(\mathsf{isList}([],c)\), while the new sentinel terminates the result.

\subsection{Append right identity: statement and induction proof}

\begin{minipage}{\linewidth}
\lstinputlisting[firstline=13,lastline=14]{theories/examples/list_rev.v}
\end{minipage}

\texttt{app\_nil\_r'} proves \(xs\mathbin{+\mkern-6mu+}[]=xs\).
Induction on \(xs\) has a reflexive empty case. In the cons case, simplification
exposes the unchanged head and an append on the tail; rewriting the induction
hypothesis completes the equality. This identity is needed when the loop exits.

\subsection{Append associativity: statement and proof}

\begin{minipage}{\linewidth}
\lstinputlisting[firstline=16,lastline=17]{theories/examples/list_rev.v}
\end{minipage}

\texttt{app\_assoc'} proves that append associates. Again, induction on the
first list makes the empty case reflexive and reduces the cons case to the
same statement for its tail. The identity changes grouping, not element order.
It is the algebraic step that lets one processed element move between the
prefix and suffix descriptions.

\subsection{Reversing append: statement and proof}

\begin{minipage}{\linewidth}
\lstinputlisting[firstline=19,lastline=24]{theories/examples/list_rev.v}
\end{minipage}

\texttt{rev\_app} proves
\(\mathsf{rev}(xs\mathbin{+\mkern-6mu+}ys)=
\mathsf{rev}(ys)\mathbin{+\mkern-6mu+}\mathsf{rev}(xs)\).
The empty case reduces to right identity of append. In the cons case,
\(\mathsf{rev}(x::xs)=\mathsf{rev}(xs)\mathbin{+\mkern-6mu+}[x]\).
The induction hypothesis moves \(\mathsf{rev}(ys)\) to the front; associativity
then reconciles the two groupings of the three appended lists.

\subsection{Reversing cons: statement and proof}

\begin{minipage}{\linewidth}
\lstinputlisting[firstline=26,lastline=27]{theories/examples/list_rev.v}
\end{minipage}

\texttt{rev\_cons} is the defining cons equation of Coq's \texttt{rev}.
Simplification already proves it. Even this small lemma is useful as a named
explanation of why prepending \(x\) to the accumulator corresponds to extending
the processed prefix at its right end.

\subsection{The invariant equation after one step}

\begin{minipage}{\linewidth}
\lstinputlisting[firstline=29,lastline=33]{theories/examples/list_rev.v}
\end{minipage}

Suppose \(v=x::v'\). Replacing the accumulator \(u\) by \(x::u\) and
remaining list \(v\) by \(v'\) preserves the recovered original list:
\[
\begin{aligned}
\mathsf{rev}(x::u)\mathbin{+\mkern-6mu+}v'
 &= (\mathsf{rev}(u)\mathbin{+\mkern-6mu+}[x])
       \mathbin{+\mkern-6mu+}v'\\
 &= \mathsf{rev}(u)\mathbin{+\mkern-6mu+}([x]\mathbin{+\mkern-6mu+}v')\\
 &= \mathsf{rev}(u)\mathbin{+\mkern-6mu+}(x::v').
\end{aligned}
\]
The proof substitutes the premise \(v=x::v'\), rewrites \texttt{rev\_cons},
uses associativity, and simplifies the singleton append. Thus the mutation
proof can delegate its pure equality to this lemma.
\subsection{The invariant equation at an empty suffix}

\begin{minipage}{\linewidth}
\lstinputlisting[firstline=35,lastline=37]{theories/examples/list_rev.v}
\end{minipage}

\texttt{rev\_prefix\_nil} records the exit identity as a direct application of
right identity. It is an explanatory corollary; the loop proof can use
\texttt{app\_nil\_r'} directly.

\subsection{Double reversal of mathematical lists: proof by induction}

\begin{minipage}{\linewidth}
\lstinputlisting[firstline=39,lastline=43]{theories/examples/list_rev.v}
\end{minipage}

\texttt{rev\_rev} proves involution:
\(\mathsf{rev}(\mathsf{rev}(xs))=xs\). The empty case is reflexive.
For a cons list, the outer reversal sees an append; \texttt{rev\_app} turns it
into the reversed singleton followed by the double-reversed tail. Simplification
and the induction hypothesis produce \(x::xs\).

\subsection{Recovering the accumulator contents from the exit equation}

\begin{minipage}{\linewidth}
\lstinputlisting[firstline=45,lastline=50]{theories/examples/list_rev.v}
\end{minipage}

\texttt{rev\_eq\_rev} derives
\(\mathsf{rev}(u)=xs\Rightarrow u=\mathsf{rev}(xs)\).
Its proof passes through \(\mathsf{rev}(\mathsf{rev}(u))\): involution gives the
first equality, and congruence under \texttt{rev} gives the second. This is what
converts the invariant equation at loop exit into the desired representation
predicate for the result.

\subsection{Opening a cons cell: definition and proof}

\begin{minipage}{\linewidth}
\lstinputlisting[firstline=108,lastline=118]{theories/examples/list_rev.v}
\end{minipage}

This lemma exposes exactly the resources in the cons clause of
\texttt{is\_list}. The pointer equality identifies the runtime value as an
object reference. The three field facts authorize reading the tag, reading the
head if needed, and reading or overwriting the link. The recursive predicate
owns the suffix independently. \texttt{iDestruct} opens the existential
witnesses and separating conjunction; \texttt{iExists}, \texttt{iSplit}, and
\texttt{iFrame} return the same resources in the stated form. No memory changes
and no ownership is duplicated by this lemma.

\subsection{Closing the new accumulator: definition and proof}

\begin{minipage}{\linewidth}
\lstinputlisting[firstline=120,lastline=128]{theories/examples/list_rev.v}
\end{minipage}

Despite its historical name, this lemma constructs a cons at the front of the
accumulator. After the current node's link has become \(a\), ownership of
\(\mathsf{isList}(u,a)\) and the three fields of that node yields
\(\mathsf{isList}(x::u,\mathsf{obj}(o))\). The proof chooses \(o,a\) as the
existential witnesses, proves the pointer equality by reflexivity, and frames
the four supplied resources. This is folding the predicate after mutation;
it is not allocation and it does not copy a node.

\subsection{One iteration: the full ownership proof}

\lstinputlisting[firstline=131,lastline=176]{theories/examples/list_rev.v}

The pure premise supplies the invariant equation with a nonempty remaining
list. The four spatial premises supply the two stack slots and both lists.
The proof opens the remaining cons cell and substitutes its pointer equality.
It then follows the source iteration in evaluation order:
\begin{enumerate}
\item \texttt{wp\_bind} with a \texttt{LetCtx} focuses on the initializer of
\texttt{nxt}. A nested \texttt{FieldLoadCtx} first loads the current pointer
from its stack slot. \texttt{wp\_field\_load} reads field \(2\), producing the
saved tail reference and returning its ownership. The stack load also returns
the stack fact. These returned facts must remain available for later commands.
\item \texttt{wp\_let} substitutes the saved reference into the continuation.
A \texttt{SeqCtx} focuses on the store while retaining both assignments as
pending code. The store's left and right evaluation contexts load \texttt{cur}
and \texttt{acc}, respectively. \texttt{wp\_field\_store} replaces only the
link fact \(o\ptso[2]\mathit{nxt}\) by \(o\ptso[2]a\).
The tag, head, old accumulator, and suffix stay framed.
\item After the store returns unit, \texttt{wp\_seq} takes the sequencing
step. The next \texttt{SeqCtx} focuses on \texttt{acc := cur}.
A \texttt{StackAssignRCtx} loads the current reference before the stack write.
\texttt{wp\_stack\_assign} replaces \(l_a\ptss a\) by
\(l_a\ptss\mathsf{obj}(o)\), retaining ownership of \(l_c\).
\item Another sequencing step reaches \texttt{cur := nxt}. The saved tail
reference is already a value. The stack write replaces
\(l_c\ptss\mathsf{obj}(o)\) by \(l_c\ptss\mathit{nxt}\).
The program now has the references claimed by the new invariant.
\item The proof chooses the saved tail as the existential witness, applies
\texttt{rev\_prefix\_step} to the pure equation, and frames the stack slots and
suffix. Finally \texttt{is\_list\_snoc\_acc} folds the modified node and old
accumulator into the new accumulator predicate.
\end{enumerate}

\texttt{iIntros "!> \_ ..."} after each primitive memory rule introduces the
later step and discards the step credit that the rule supplies. The remaining
named hypothesis is the returned ownership. Discarding a credit is harmless;
discarding the link or stack ownership would prevent a later access.
The code uses explicit evaluation contexts because a store or assignment
returns a value inside a pending sequence: the residual sequencing step must
be handled before proving the next command. The bind rule arranges the
continuation; it does not itself execute that sequencing step.

\subsection{The exit assertion: \texttt{rev\_done}}

\begin{minipage}{\linewidth}
\lstinputlisting[firstline=179,lastline=181]{theories/examples/list_rev.v}
\end{minipage}

An invariant alone does not imply that the loop has finished. For example, the
entry state satisfies \texttt{rev\_inv} with the entire input still unprocessed.
The loop's postcondition therefore separately states that \texttt{cur} owns an
empty list and \texttt{acc} owns the reversed input. Both stack facts are
retained for the enclosing variable scopes. This is a stronger assertion than
merely re-establishing the invariant, and avoids an invalid inference that
\(\mathsf{rev}(u)\mathbin{+\mkern-6mu+}v=xs\) forces \(v=[]\).

\subsection{The loop theorem: induction and both branches}

\lstinputlisting[firstline=183,lastline=225]{theories/examples/list_rev.v}

The induction is on the mathematical remaining list \(v\). The statement
quantifies over \(u,a,c\) after fixing \(v\), so the induction hypothesis is
available for any new accumulator and current reference. Fixing those values
before induction would make the hypothesis too weak for the updates.

\paragraph{Empty remaining list.}
Opening \(\mathsf{isList}([],c)\) identifies the sentinel and gives its false
tag. \texttt{wp\_while} unfolds one loop test into an \texttt{if}.
The \texttt{IfCtx} retains the two branches while the nested load context
reads the current reference and its tag. The field load returns false and
\texttt{wp\_if\_false} selects unit. There is no body iteration in this case.
Right identity reduces the invariant equation to \(\mathsf{rev}(u)=xs\);
\texttt{rev\_eq\_rev} changes it to \(u=\mathsf{rev}(xs)\).
The proof supplies both references to \texttt{rev\_done}, frames the accumulator
and stack facts, and reconstructs the empty sentinel predicate from all three
fields. The guard consumed none of the sentinel's ownership.

\paragraph{Nonempty remaining list.}
Opening \(\mathsf{isList}(x::v,c)\) gives the true tag and saved-tail witness.
The guard proof reads the tag and selects the true branch. To invoke the body
lemma, it must reassemble the cons predicate from the returned tag, head,
link, and suffix ownership. A \texttt{SeqCtx} isolates the iteration from the
recursive loop. \texttt{wp\_wand} first proves the iteration with its own
postcondition, then supplies that postcondition to the remaining proof.

The resulting existential tail reference and pure equation describe
\(x::u\) and \(v\). After \texttt{wp\_seq}, the induction hypothesis applies to
this strictly smaller remaining list with the updated references. It produces
\texttt{rev\_done} directly. The proof never assumes that the loop has exited
while a cons remains, and never needs to guess a final accumulator from the
original invariant equation alone.

\subsection{Packaging the invariant: \texttt{wp\_rev\_loop}}

\begin{minipage}{\linewidth}
\lstinputlisting[firstline=227,lastline=234]{theories/examples/list_rev.v}
\end{minipage}

This wrapper opens the four existential witnesses of \texttt{rev\_inv} and
its pure equality, then applies \texttt{wp\_rev\_loop\_run} with those witnesses
and all four spatial resources. There is no additional operational argument.
It is the interface used by the function proof, which can initialize an
invariant without exposing its witnesses at the call site.

\subsection{The complete function proof}

\lstinputlisting[firstline=236,lastline=266]{theories/examples/list_rev.v}

First \texttt{is\_list\_obj} identifies the input value as an object reference
while retaining its list ownership. This matters because the language's value
type also includes closures: substituting a variable through an arbitrary
value need not leave that value syntactically unchanged. An object literal has
no free variables, so the substitutions reduce as required.

The \texttt{AppLCtx} evaluates the closure, \texttt{wp\_rec} creates it, and
\texttt{wp\_app} substitutes the object argument into the body. A
\texttt{LetCtx} focuses on \texttt{call\_nil}. The proof uses
\texttt{wp\_wand} with an empty resource selection to invoke \texttt{nil\_spec}:
the allocator requires no input ownership, and the original list remains
framed outside this call. The returned sentinel gives \(\mathsf{isList}([],a_0)\).

The two \texttt{wp\_var\_bind} applications allocate the accumulator and current
stack slots. Their postconditions require returning ownership of each slot at
scope exit. After substitution the loop is \texttt{rev\_loop\_at}, and the last
expression is \texttt{sld} of the accumulator slot. A \texttt{SeqCtx} focuses
on the loop. Its invariant uses witnesses \([],xs,a_0,p\); the pure equality
simplifies to \(xs=xs\), and the two stack facts and two list predicates are
exactly the owned resources already available.

\texttt{wp\_rev\_loop} yields the stronger exit assertion.
After the final sequencing step, \texttt{wp\_stack\_load} returns the accumulator
reference and its slot ownership. The two \texttt{iExists} and \texttt{iFrame}
steps satisfy the nested variable scopes' requirements in the reverse order
of their introduction: current slot first, accumulator slot second.
The result predicate \(\mathsf{isList}(\mathsf{rev}(xs),a)\) remains for the
client. Ownership of the old sentinel may be discarded because Iris is
affine; this is not a deallocation step in the machine.

\subsection{A client equation: reversing twice}

\begin{minipage}{\linewidth}
\lstinputlisting[firstline=268,lastline=270]{theories/examples/list_rev.v}
\end{minipage}

This client composes two calls using an immutable intermediate reference.
\subsection{The double-reversal client theorem and its proof}

\begin{minipage}{\linewidth}
\lstinputlisting[firstline=273,lastline=285]{theories/examples/list_rev.v}
\end{minipage}

The proof focuses on the first call with \texttt{LetCtx}, applies
\texttt{reverse\_spec}, and receives ownership of a list with contents
\(\mathsf{rev}(xs)\). After substituting the returned reference, it applies the
same theorem a second time. \texttt{rev\_rev} rewrites the second result's
contents to \(xs\). This is an equation of represented sequences.
Each call creates a fresh empty sentinel, so it is not a theorem that every
heap field or the terminal sentinel's identity is restored.

\subsection{The printed Cangjie program}

\lstinputlisting{generated/cj/list_rev.cj}

\texttt{p\_list\_rev} in \texttt{emit\_cj.v} declares the same \texttt{Node}
layout, reuses the five list constructors and observers, and adds
\texttt{reverse} with \texttt{rev\_body}. Its closed example builds a singleton
and reverses it. The printer renders the source \texttt{VarBind}s as mutable
Cangjie locals and field \(2\) as \texttt{next}. The generated file is standalone
so its node declarations do not collide with \texttt{list.cj} in the tests.
It is included in the normal \texttt{make cj} pipeline and the expected-output
check. The executable under test is this generated file, without a separately
handwritten reversal implementation.

\subsection{Executable tests and what they check}

\lstinputlisting{test/list_rev_test.cj}

The six test cases cover empty input, a singleton, four distinct elements,
duplicates and negative values, double reversal, and mutation through saved
references while a separate list remains unchanged. The nonempty sequence
tests inspect every expected element and check the terminal empty tag, so a
missing or extra node is observable. The mutation test saves references to all
three original cons nodes and observes their changed links. A copying
implementation that left those original links unchanged would fail that test.
The second list tests the concrete behavior expected from framing.

These are finite runtime checks of the printed Cangjie on the installed
compiler. They do not replace the general theorem over arbitrary mathematical
lists. Conversely the Coq proof checks the CJR term and does not prove the
compiler correct. The test command is
\begin{lstlisting}
cjc generated/cj/list_rev.cj test/list_rev_test.cj \
    --test -o test/bin/list_rev_test
test/bin/list_rev_test
\end{lstlisting}
The dated test report records the actual run. The assumption audit includes
\texttt{reverse\_spec} and \texttt{reverse\_twice\_spec}; neither theorem is
closed using an admitted obligation or an added axiom.


\section{Arrays}
\label{sec:arrays}

The file \texttt{theories/examples/array.v} is a library, in the same sense as the list: the language has no array type, and the library builds one from a final class and a raw block.
An array of \(n\) integers is an object with two fields and a contiguous block of \(n\) words.
Field \(0\) holds \(n\).
Field \(1\) holds the base pointer of the block.
Word \(i\) of the block holds element \(i\).
The index is a computed integer, so the elements cannot be object fields: a field index in the core is a static natural number, whereas \texttt{Offset} and \texttt{Load} take a value.

\texttt{Alloc} of zero words is stuck, so \texttt{make} asks for a positive length.
An index outside \(0,\ldots,n-1\) is likewise stuck: the specification owns only the words of the block, and a load with no points-to fact does not reduce.

\subsection{The library}

In Cangjie the value is a non-\texttt{open} class.
Allocation of \(n\) words is \texttt{malloc} of \(n\) times the size of \texttt{Int64}; the core step is \texttt{Alloc} of the word count \(n\), which returns a pointer to \(n\) zero words.
\texttt{get} and \texttt{set} take the array and the index, and \texttt{set} also takes the new element.
Core application is unary, so those calls lower to one struct argument, as \texttt{cons} did.

\lstinputlisting{generated/cj/array.cj}

The core terms are the same operations with the struct lowering made explicit.

\begin{lstlisting}
Definition make_body : expr :=
  Let (Some "p") (Alloc (Var "n"))
    (New [] [Var "n"; Var "p"]).

Definition get_body : expr :=
  Let (Some "arr") (StructLoad (Var "args") 0)
    (Let (Some "i") (StructLoad (Var "args") 1)
      (Let (Some "p") (FieldLoad (Var "arr") 1)
        (Load (Offset (Var "p") (Var "i"))))).

Definition set_body : expr :=
  Let (Some "arr") (StructLoad (Var "args") 0)
    (Let (Some "i") (StructLoad (Var "args") 1)
      (Let (Some "x") (StructLoad (Var "args") 2)
        (Let (Some "p") (FieldLoad (Var "arr") 1)
          (Store (Offset (Var "p") (Var "i")) (Var "x"))))).
\end{lstlisting}

\texttt{length} is a field load of index \(0\).

\subsection{The representation}

\begin{definition}[Array representation]
\(\mathsf{isArray}(xs, a)\) holds when \(a = \mathsf{obj}(o)\), the base pointer is \(p\), and
\[
\begin{aligned}
&o \ptso[0] \lvert xs \rvert
\sep
o \ptso[1] \mathsf{ptr}(p)
\sep
p \ptsb \lvert xs \rvert
\sep
\mathop{\ast}\limits_{i < \lvert xs \rvert}
  (p+i) \ptsr xs_i.
\end{aligned}
\]
\end{definition}

The block fact is the right to treat \([p, p+\lvert xs \rvert)\) as one allocation.
The separating conjunction of the words is what lets a proof load or store one index and still own the others.
Two different indices are two different addresses, so the facts do not overlap.

The Coq predicate is that assertion.
\texttt{array\_cells} is the big separating conjunction, indexed by position in \(xs\).

\begin{lstlisting}
Definition array_cells (base : Z) (xs : list Z) : iProp S :=
  (* separating conjunction over i |-> xs_i at address base+i *)
  [ * list] i |-> x in xs,
    RawId (base + Z.of_nat i) |->r LitV (LitInt x).

Definition is_array (xs : list Z) (a : val) : iProp S :=
  exists o base, pure (a = LitV (LitObj o)) *
    ObjId o |->o[0] LitV (LitInt (Z.of_nat (length xs))) *
    ObjId o |->o[1] LitV (LitPtr base) *
    RawId base |->b (length xs) *
    array_cells base xs.
\end{lstlisting}

\subsection{Allocation of \(n\) words}

The one-word rule is not enough for \texttt{make}.
The machine already steps \texttt{Alloc} of any positive word count: it zeroes that many words, records the length at the base, and advances the counter by the count.
\texttt{wp\_alloc} is the matching logic rule.
From \(0 < n\),
\[
\Wpre\,\mathsf{alloc}(n)
\,\Bigl\{\,v.\;
  \exists \ell.\;
  v = \mathsf{ptr}(\ell)
  \sep
  \ell \ptsb n
  \sep
  \mathop{\ast}\limits_{i < n} (\ell+i) \ptsr 0
\,\Bigr\}.
\]
The proof is the one-word proof with a finite map of cells instead of a singleton.
\texttt{init\_cells} builds that map, \texttt{wf\_alloc} shows the resulting machine is well formed, and the ghost-map update hands back one points-to fact per word.

\subsection{What each operation establishes}

\paragraph{Make.}
For \(0 < n\),
\[
\Wpre\,\mathsf{make}(n)
\,\{\,a.\; \mathsf{isArray}(\mathsf{replicate}(n, 0), a)\,\}.
\]
Allocation produces the zero words and the block fact.
The following \texttt{new} stores \(n\) and the pointer in the object.
This is \texttt{make\_spec}.

\paragraph{Length.}
\[
\mathsf{isArray}(xs, a)
\vdash
\Wpre\,\mathsf{length}(a)
\,\{\,v.\; \pure{v = \lvert xs \rvert} \sep \mathsf{isArray}(xs, a)\,\}.
\]
The proof loads field \(0\) and returns the same predicate.
Length is an observation.
This is \texttt{length\_spec}.

\paragraph{Get.}
If \(xs_i = x\), writing \(i\) for the mathematical index,
\[
\mathsf{isArray}(xs, a)
\vdash
\Wpre\,\mathsf{get}(a, i)
\,\{\,v.\; \pure{v = x} \sep \mathsf{isArray}(xs, a)\,\}.
\]
The body reads the pointer from field \(1\), adds \(i\), and loads that word.
The lookup lemma for a big separating conjunction lends out the one cell and takes it back, so the other cells are not mentioned by the load rule and stay in the predicate.
This is \texttt{get\_spec}.

\paragraph{Set.}
If index \(i\) is in range,
\[
\mathsf{isArray}(xs, a)
\vdash
\Wpre\,\mathsf{set}(a, i, x)
\,\{\,v.\; \pure{v = \Unit} \sep \mathsf{isArray}(xs[i \mapsto x], a)\,\}.
\]
The store updates that one word.
The length stored in the object and the block length are unchanged, because replacing an element does not change \(\lvert xs \rvert\).
This is \texttt{set\_spec}.

\subsection{The array equations}

The programs are the compositions of the functions above.
Each one is a \texttt{let} of one call followed by the next, not a fused body.

\begin{theorem}[Read after write]
If index \(i\) is in range, then
\[
\mathsf{isArray}(xs, a)
\vdash
\Wpre\,(\mathsf{get}(\mathsf{set}(a, i, x), i))
\,\{\,v.\; \pure{v = x} \sep \mathsf{isArray}(xs[i \mapsto x], a)\,\}.
\]
If \(j\) is also in range and \(i \neq j\), then
\[
\mathsf{isArray}(xs, a)
\vdash
\Wpre\,(\mathsf{get}(\mathsf{set}(a, i, x), j))
\,\{\,v.\; \pure{v = xs_j} \sep \mathsf{isArray}(xs[i \mapsto x], a)\,\}.
\]
\end{theorem}

These are \texttt{set\_then\_get\_same} and \texttt{set\_then\_get\_other}.
The first is \texttt{set\_spec} followed by \texttt{get\_spec} at the updated list, whose lookup at \(i\) is \(x\).
The second uses the same \texttt{set\_spec}, then \texttt{get\_spec} at \(j\).
The lookup of an updated list at a different index is the old element.
The separating conjunction is why the store at \(i\) does not have a points-to fact for \(j\) among its hypotheses: that fact is framed, so it cannot change.

\begin{theorem}[Length is invariant]
If index \(i\) is in range, then
\[
\mathsf{isArray}(xs, a)
\vdash
\Wpre\,(\mathsf{length}(\mathsf{set}(a, i, x)))
\,\{\,v.\; \pure{v = \lvert xs \rvert} \sep \mathsf{isArray}(xs[i \mapsto x], a)\,\}.
\]
For \(0 < n\) and \(i < n\),
\[
\begin{aligned}
\Wpre\,(\mathsf{length}(\mathsf{make}(n)))
&\,\{\,v.\; \exists a.\; \pure{v = n} \sep \mathsf{isArray}(\mathsf{replicate}(n, 0), a)\,\},\\
\Wpre\,(\mathsf{get}(\mathsf{make}(n), i))
&\,\{\,v.\; \exists a.\; \pure{v = 0} \sep \mathsf{isArray}(\mathsf{replicate}(n, 0), a)\,\}.
\end{aligned}
\]
\end{theorem}

The length equation is \texttt{set\_then\_length\_spec}: after \texttt{set}, \texttt{length} reads the field that still holds \(\lvert xs \rvert\).
The two \texttt{make} equations are \texttt{make\_then\_length\_spec} and \texttt{make\_then\_get\_spec}.
A fresh array is a vector of zeroes of the requested length, and reading it does not change that fact.

What this development does not prove is an arbitrary loop over the array, such as a sum of every cell.
The cells are already a separating conjunction indexed by the mathematical list, which is the induction hypothesis such a loop would use.
The equations above are the base: read and write agree with list update, and they do not disturb the length or the other cells.

\section{Array lists}
\label{sec:array-lists}

The file \texttt{theories/examples/array\_list.v} lowers the resizable-buffer core of Cangjie's \texttt{ArrayList}: a default constructor, the live length, the capacity, emptiness, reading and writing an index, append, removal at an index, growth, and the version counter.
The source class continues past that core, with an iterator, the \texttt{Collection} interface, ranges, higher-order methods, and an embedded sort.
Those stay out.
The core is monomorphic in \texttt{Int64}, and it uses the array library already proved: a final object, a raw block, and a word-by-word copy.

\texttt{Alloc} of zero words is stuck, so the only constructor is the default one, at capacity \(16\).
There is no shift, so growth doubles by adding the capacity to itself.
There are no exceptions.
An index outside \([0, \mathit{size})\) lowers to a load of a non-pointer, which has no step, and every specification takes the arm that is in range.
\texttt{set} leaves the version alone.
\texttt{add} and \texttt{remove} increment it.
Growth replaces the buffer and drops the old block and the old array object.
The source relies on collection for that drop.
This language has a one-word free and no collection, so the old facts are not returned to the allocator.
Iris is affine, and the proof simply does not use them again.

\subsection{The library}

In Cangjie the value is a non-\texttt{open} class.
Field \(0\) is the buffer, itself the array of the previous section.
Field \(1\) is the live length.
Field \(2\) is the version.
The buffer's own field \(0\) is the capacity, and its field \(1\) is the base pointer of a block of that many words.
The live prefix of the block is the mathematical list.
The tail is zero, which is what \texttt{Alloc} writes and what removal writes back.

\lstinputlisting{generated/cj/array_list.cj}

Core application is unary, so each method lowers to a closure applied to one struct of arguments, as \texttt{get} and \texttt{set} did for arrays.
The pure model of the three mutating operations is ordinary list algebra.

\begin{lstlisting}
Definition al_set (xs : list Z) (i : nat) (x : Z) : list Z := <[i := x]> xs.
Definition al_add (xs : list Z) (x : Z) : list Z := xs ++ [x].
Definition al_remove (xs : list Z) (i : nat) : list Z :=
  take i xs ++ drop (S i) xs.
\end{lstlisting}

Appending puts \(x\) at index \(\lvert xs \rvert\) and leaves every earlier index unchanged.
Updating index \(i\) reads back \(x\) there and leaves every other index unchanged.
Removing index \(i\) shortens the list by one and shifts the suffix down.

\subsection{The representation}

Two objects.
The inner one is the buffer.
The outer one is the list.

\begin{definition}[Array-list representation]
Write \(\mathsf{slots}(xs, c)\) for \(xs\) followed by \(c - \lvert xs \rvert\) zeroes, which is defined only for \(\lvert xs \rvert \le c\).
\(\mathsf{isBuf}(\mathit{slots}, b)\) is the array predicate of the previous section with the whole capacity in place of the live prefix: field \(0\) holds \(\lvert \mathit{slots} \rvert\), field \(1\) holds the base pointer, and the block owns every word of \(\mathit{slots}\).
\(\mathsf{isArrayList}(xs, c, v, a)\) holds when \(0 < c\), \(\lvert xs \rvert \le c\), \(a = \mathsf{obj}(o)\), and
\[
\begin{aligned}
&o \ptso[0] b
\sep
o \ptso[1] \lvert xs \rvert
\sep
o \ptso[2] v
\sep
\mathsf{isBuf}(\mathsf{slots}(xs, c), b).
\end{aligned}
\]
\end{definition}

The Coq predicates are \texttt{slots\_of}, \texttt{is\_buf}, and \texttt{is\_arraylist}.
The cells are \texttt{array\_cells}, reused from the array library, so a proof of a single load or store is the same lookup lemma as \texttt{get} and \texttt{set} there.
The zero tail is part of the owned block.
Appending into spare capacity overwrites the first zero.
Removing the last live word writes a zero back, so the tail stays zero.

\subsection{Copy and growth}

Copy is one loop.
It binds an index at \(0\) and, while that index plus one is at most the length, loads one word and stores it one step further along.
The loop is \texttt{copy\_vars}.
Two invariants share that body.

The first is a copy between two blocks, used by growth.
Source and destination do not overlap, so each iteration may read the source and write the destination in either order.
\texttt{copy\_disjoint} says that after \(\mathit{len}\) steps the destination holds the source chunk at the destination offset, and the source is unchanged.

The second is a forward shift on one block, used by removal.
The source offset is the destination offset plus one, so each word is read before the next store overwrites it.
\texttt{copy\_shift} says that the live prefix is the list with the suffix shifted down by one.
A left shift of length zero, which is removal of the last element, is the loop condition failing at once.

Growth asks for a minimum capacity.
The new capacity is \(2c\) when the minimum is at most \(2c\), and the minimum otherwise.
\[
\mathsf{growCap}(c, m) =
\begin{cases}
2c & \text{if } m \le 2c,\\
m & \text{otherwise.}
\end{cases}
\]
Both arms allocate a positive block, copy the live prefix into it, and store the new array object in field \(0\).
The mathematical list and the version are unchanged.
The old block remains allocated and is not mentioned by the postcondition.
This is \texttt{grow\_spec}.
The Coq name of the capacity is \texttt{grow\_cap}; doubling is \texttt{doubled}, which is addition because the core has no shift.

Append uses that rule only when the live length already equals the capacity.
Then the minimum is the length plus one, which is at most \(2c\) for \(c > 0\), so the capacity doubles.
When the length is strictly below the capacity, append writes into the zero tail and the capacity stays.

\subsection{What each operation establishes}

\paragraph{Init.}
\[
\Wpre\,\mathsf{init}()
\,\{\,a.\; \mathsf{isArrayList}([], 16, 0, a)\,\}.
\]
Allocation of \(16\) words produces the zero block.
The inner object stores \(16\) and the pointer.
The outer object stores that buffer, length \(0\), and version \(0\).
This is \texttt{init\_spec}.
\texttt{default\_capacity} is \(16\).

\paragraph{Size, capacity, emptiness, version.}
Each is a field load, or a load followed by a comparison with zero.
\[
\begin{aligned}
\mathsf{isArrayList}(xs, c, v, a)
&\vdash
\Wpre\,\mathsf{size}(a)
\,\{\,r.\; \pure{r = \lvert xs \rvert} \sep \mathsf{isArrayList}(xs, c, v, a)\,\},\\
\mathsf{isArrayList}(xs, c, v, a)
&\vdash
\Wpre\,\mathsf{capacity}(a)
\,\{\,r.\; \pure{r = c} \sep \mathsf{isArrayList}(xs, c, v, a)\,\},\\
\mathsf{isArrayList}(xs, c, v, a)
&\vdash
\Wpre\,\mathsf{isEmpty}(a)
\,\{\,r.\; \pure{r = (\lvert xs \rvert = 0)} \sep \mathsf{isArrayList}(xs, c, v, a)\,\},\\
\mathsf{isArrayList}(xs, c, v, a)
&\vdash
\Wpre\,\mathsf{version}(a)
\,\{\,r.\; \pure{r = v} \sep \mathsf{isArrayList}(xs, c, v, a)\,\}.
\end{aligned}
\]
These are \texttt{size\_spec}, \texttt{capacity\_spec}, \texttt{is\_empty\_spec}, and \texttt{version\_spec}.
Capacity reads the buffer's length field, which \texttt{slots\_of} keeps equal to \(c\).

\paragraph{Get and set.}
If \(xs_i = y\),
\[
\begin{aligned}
\mathsf{isArrayList}(xs, c, v, a)
&\vdash
\Wpre\,\mathsf{get}(a, i)
\,\{\,r.\; \pure{r = y} \sep \mathsf{isArrayList}(xs, c, v, a)\,\},\\
\mathsf{isArrayList}(xs, c, v, a)
&\vdash
\Wpre\,\mathsf{set}(a, i, x)
\,\{\,r.\; \pure{r = \Unit} \sep \mathsf{isArrayList}(xs[i \mapsto x], c, v, a)\,\}.
\end{aligned}
\]
Both go through the buffer pointer and the raw word at index \(i\).
The range test is two comparisons, \(c\) against the index and the index against \(-1\), and the specification is the success arm.
\texttt{set} stores into the live prefix and does not touch field \(2\).
These are \texttt{al\_get\_spec} and \texttt{al\_set\_spec}.
The names are prefixed because the array library already exports \texttt{get\_spec} and \texttt{set\_spec}.

\paragraph{Add.}
\[
\mathsf{isArrayList}(xs, c, v, a)
\vdash
\Wpre\,\mathsf{add}(a, x)
\,\{\,r.\; \pure{r = \Unit} \sep
  \mathsf{isArrayList}(xs \mathbin{+\mkern-6mu+} [x],\, \mathsf{addCap}(c, \lvert xs \rvert),\, v+1, a)\,\}.
\]
\(\mathsf{addCap}(c, n)\) is \(c\) when \(n < c\), and \(\mathsf{growCap}(c, n+1)\) otherwise.
The body writes \(x\) at the old length, stores the new length, and increments the version.
When the buffer is already full it first calls \texttt{grow}.
This is \texttt{add\_spec}.

\paragraph{Remove.}
If \(xs_j = y\),
\[
\mathsf{isArrayList}(xs, c, v, a)
\vdash
\Wpre\,\mathsf{remove}(a, j)
\,\{\,r.\; \pure{r = y} \sep
  \mathsf{isArrayList}(\mathsf{delete}(xs, j),\, c,\, v+1, a)\,\}.
\]
\(\mathsf{delete}(xs, j)\) is the prefix of length \(j\) followed by the suffix after \(j\).
The body loads \(y\) before the shift, copies the suffix down by one, writes zero at the old last index, stores the new length, and increments the version.
The capacity is the same block.
This is \texttt{remove\_spec}.

\subsection{The array-list equations}

The programs are compositions of the calls above.
Each one is a \texttt{let} of one call followed by the next, proved from the triples by framing, as \texttt{set\_then\_get} was for arrays.

\begin{theorem}[A fresh list]
\[
\begin{aligned}
\Wpre\,(\mathsf{size}(\mathsf{init}()))
&\,\{\,r.\; \exists a.\; \pure{r = 0} \sep \mathsf{isArrayList}([], 16, 0, a)\,\},\\
\Wpre\,(\mathsf{isEmpty}(\mathsf{init}()))
&\,\{\,r.\; \exists a.\; \pure{r = \True} \sep \mathsf{isArrayList}([], 16, 0, a)\,\},\\
\Wpre\,(\mathsf{capacity}(\mathsf{init}()))
&\,\{\,r.\; \exists a.\; \pure{r = 16} \sep \mathsf{isArrayList}([], 16, 0, a)\,\},\\
\Wpre\,(\mathsf{version}(\mathsf{init}()))
&\,\{\,r.\; \exists a.\; \pure{r = 0} \sep \mathsf{isArrayList}([], 16, 0, a)\,\}.
\end{aligned}
\]
\end{theorem}

Size and emptiness are \texttt{init\_then\_size\_spec} and \texttt{init\_then\_empty\_spec}.
Capacity and version are \texttt{init\_then\_capacity\_spec} and \texttt{init\_then\_version\_spec}.

\begin{theorem}[Append]
\[
\begin{aligned}
&\mathsf{isArrayList}(xs, c, v, a) \vdash
  \Wpre\,(\mathsf{size}(\mathsf{add}(a, x)))\\
&\qquad \{\,r.\; \pure{r = \lvert xs \rvert + 1} \sep
  \mathsf{isArrayList}(xs \mathbin{+\mkern-6mu+} [x],\,
    \mathsf{addCap}(c, \lvert xs \rvert),\, v+1, a)\,\},\\
&\mathsf{isArrayList}(xs, c, v, a) \vdash
  \Wpre\,(\mathsf{get}(\mathsf{add}(a, x), \lvert xs \rvert))\\
&\qquad \{\,r.\; \pure{r = x} \sep
  \mathsf{isArrayList}(xs \mathbin{+\mkern-6mu+} [x],\,
    \mathsf{addCap}(c, \lvert xs \rvert),\, v+1, a)\,\}.
\end{aligned}
\]
If \(i < \lvert xs \rvert\) and \(xs_i = y\), the same append followed by \texttt{get} at \(i\) returns \(y\).
The capacity is \(c\) when \(\lvert xs \rvert < c\), and \(2c\) when \(\lvert xs \rvert = c\).
The version after append is \(v+1\).
\end{theorem}

Size and the new-index lookup are \texttt{add\_then\_size\_spec} and \texttt{add\_then\_get\_new\_spec}.
An older index is \texttt{add\_then\_get\_old\_spec}.
Spare capacity is \texttt{add\_then\_capacity\_room\_spec}, a full buffer is \texttt{add\_then\_capacity\_full\_spec}, and the version is \texttt{add\_then\_version\_spec}.
The lookup at the new index is the lemma that an appended element sits at the old length.
The lookup at an older index is the lemma that append does not change a shorter prefix, including across the growth that replaces the block.

\begin{theorem}[Update]
If \(xs_i = y\), then reading index \(i\) after writing \(x\) there returns \(x\), and the version is still \(v\).
If \(j\) is also in range and \(i \neq j\), reading \(j\) returns \(xs_j\).
The length after the update is still \(\lvert xs \rvert\).
\end{theorem}

The two lookups are \texttt{set\_then\_get\_same\_spec} and \texttt{set\_then\_get\_other\_spec}.
Length and version are \texttt{set\_then\_size\_spec} and \texttt{set\_then\_version\_spec}.
The version equation is the observation that \texttt{set} does not store field \(2\).

\begin{theorem}[Append then remove]
\[
\begin{aligned}
&\mathsf{isArrayList}(xs, c, v, a) \vdash
  \Wpre\,(\mathsf{remove}(\mathsf{add}(a, x), \lvert xs \rvert))\\
&\qquad \{\,r.\; \pure{r = x} \sep
  \mathsf{isArrayList}(xs,\, \mathsf{addCap}(c, \lvert xs \rvert),\, v+2, a)\,\}.
\end{aligned}
\]
Removing an in-range index and then reading the size returns \(\lvert xs \rvert - 1\).
\end{theorem}

The round trip is \texttt{add\_then\_remove\_last\_spec}.
Removal of the element just appended restores \(xs\).
The capacity is whatever append left behind, which may be \(2c\), and the version has been incremented twice.
The size equation is \texttt{remove\_then\_size\_spec}.

What this development does not prove is the iterator, equality of two lists, or a bulk operation such as \texttt{filter}.
The predicate already owns the live prefix and the zero tail as one separating conjunction, which is the induction hypothesis a further loop would use.
The equations above are the base: length, lookup, update, append, and removal agree with the list algebra, growth only changes the capacity, and the version moves only on append and removal.

\input{generated/cj/further.tex}

\section{Quicksort}
\label{sec:qsort}

\texttt{theories/examples/qsort.v} is an in-place Lomuto quicksort of \texttt{Int64}.
One Gallina function is the functional specification.
One core-language program, parameterised by length, read, and write, is proved to compute that function.
The fixed array and the array list each supply those three operations.
The algorithm applies their specifications and leaves \(\mathsf{isArray}\) and \(\mathsf{isArrayList}\) folded.

The slice is half-open, \([lo, hi)\).
The pivot is the element at index \(hi - 1\).
A scan walks \(j\) from \(lo\) up to that index.
Each value \(\le\) the pivot is swapped into the stored prefix and the prefix boundary \(i\) advances.
The pivot is then swapped to \(i\).
The recursive calls sort \([lo, i)\) and \([i+1, hi)\).
The exported operation reads the length \(n\) and sorts \([0, n)\): every word of an array, or the live prefix of an array list.
The arithmetic in the program is subtraction and \(\le\).

\subsection{The functional algorithm}

Lists are mathematical lists of integers.
\(xs_k\) means the lookup \(xs !! k\), and \(xs[k \mapsto z]\) means the update \(\langle k := z \rangle\, xs\).
A missing lookup leaves the list alone.

\begin{lstlisting}
Definition swap_list (xs : list Z) (i j : nat) : list Z :=
  match xs !! i, xs !! j with
  | Some xi, Some xj => <[i := xj]> (<[j := xi]> xs)
  | _, _ => xs
  end.
\end{lstlisting}

\(\mathsf{swap}(xs, i, j)\) exchanges the elements at two in-range indexes.
The code performs the two writes in the other order, first index \(i\) and then index \(j\).
When both lookups exist, the two orders agree: if \(i = j\) both are the identity, and if \(i \neq j\) the updates commute.
That equation is \texttt{swap\_list\_code}.

\begin{lstlisting}
Lemma swap_list_code xs i j xi xj :
  xs !! i = Some xi -> xs !! j = Some xj ->
  <[j := xi]> (<[i := xj]> xs) = swap_list xs i j.
\end{lstlisting}

The scan is a fuelled recursion.
The fuel \(n\) is the number of indexes still to visit.
\(i\) is the boundary of the stored prefix, \(j\) is the next index to read, and \(\mathit{bound}\) is the index of the pivot, which the scan does not move.

\begin{lstlisting}
Fixpoint lomuto (n : nat) (xs : list Z) (pivot : Z) (i j bound : nat)
  : list Z * nat :=
  match n with
  | O => (xs, i)
  | S n =>
      if Nat.ltb j bound then
        match xs !! j with
        | Some v =>
            if Z.leb v pivot then
              lomuto n (swap_list xs i j) pivot (S i) (S j) bound
            else lomuto n xs pivot i (S j) bound
        | None => (xs, i)
        end
      else (xs, i)
  end.
\end{lstlisting}

While \(j < \mathit{bound}\) and index \(j\) holds \(v\), a value \(v \le \mathit{pivot}\) is exchanged with index \(i\) and both cursors advance.
A value \(v > \mathit{pivot}\) leaves the list unchanged and advances only \(j\).
The result is the list and the final boundary.

Partition reads the pivot, runs the scan from \(i = j = lo\) with fuel \((hi - 1) - lo\), and exchanges the pivot with the returned boundary.

\begin{lstlisting}
Definition partition_list (xs : list Z) (lo hi : nat) : list Z * nat :=
  match xs !! (hi - 1)%nat with
  | Some pivot =>
      match lomuto ((hi - 1) - lo)%nat xs pivot lo lo (hi - 1)%nat with
      | (zs, q) => (swap_list zs q (hi - 1)%nat, q)
      end
  | None => (xs, lo)
  end.
\end{lstlisting}

The recursive sort is also fuelled.
The fuel is at least the length of the slice.
A slice of length at most one is already finished.
A longer slice is partitioned, the left part is sorted, and the right part is sorted on that result.
Both recursive calls reuse the predecessor fuel: each of those slices has length at most \(hi - lo - 1\).

\begin{lstlisting}
Fixpoint qsort_slice (n : nat) (xs : list Z) (lo hi : nat) : list Z :=
  match n with
  | O => xs
  | S n =>
      if Nat.leb hi (lo + 1)%nat then xs
      else
        let '(ys, p) := partition_list xs lo hi in
        qsort_slice n (qsort_slice n ys lo p) (S p) hi
  end.

Definition qsort (xs : list Z) : list Z :=
  qsort_slice (length xs) xs 0%nat (length xs).
\end{lstlisting}

\(\mathsf{qsort}(xs)\) sorts the whole list.
The test \(\mathit{hi} \le lo + 1\) is the Gallina form of the program test \((lo + 1) \le (hi - 1)\) failing.
On a slice of length \(0\) or \(1\) the function returns its input.

\subsection{Sortedness and permutation}

A slice is sorted when every pair of elements in it is in non-decreasing order.
Indexes outside the slice are not constrained.

\begin{lstlisting}
Definition slice_sorted (xs : list Z) (lo hi : nat) : Prop :=
  forall i j x y, (lo <= i)%nat -> (i < j)%nat -> (j < hi)%nat ->
    xs !! i = Some x -> xs !! j = Some y -> (x <= y)%Z.
\end{lstlisting}

In mathematics, \(\mathsf{sorted}(xs, lo, hi)\) means
\[
lo \le i < j < hi \land xs_i = x \land xs_j = y \implies x \le y.
\]
Equal values satisfy the comparison in either order.
The specification records the multiset of values, which a permutation preserves, and this order.

The slice specification is the induction hypothesis of the fuel recursion.
\[
ys = \mathsf{qsortSlice}(n, xs, lo, hi).
\]

The source writes a permutation as \(x \equiv_p y\), a pure assertion in corner brackets, and a stack cell as \(l \mapsto_s v\).
The listings spell those with \texttt{Permutation}, \texttt{bi\_pure}, and \texttt{stack\_mapsto}.
The wand is written \texttt{-*}.

\begin{lemma}[Slice]
If \(hi \le \lvert xs \rvert\) and \(hi - lo \le n\), then \(\lvert ys \rvert = \lvert xs \rvert\), \(ys\) is a permutation of \(xs\), \(ys\) and \(xs\) agree at every index outside \([lo, hi)\), and \(\mathsf{sorted}(ys, lo, hi)\).
\end{lemma}

\begin{lstlisting}
Lemma qsort_slice_spec n xs lo hi :
  (hi <= length xs)%nat ->
  (hi - lo <= n)%nat ->
  let ys := qsort_slice n xs lo hi in
  length ys = length xs /\ Permutation ys xs /\
  (forall k, (k < lo \/ hi <= k)%nat -> ys !! k = xs !! k) /\
  slice_sorted ys lo hi.
\end{lstlisting}

The whole-list facts are that lemma at \(lo = 0\), \(hi = \lvert xs \rvert\), and \(n = \lvert xs \rvert\).

\begin{lemma}[Permutation]
\(\mathsf{qsort}(xs)\) is a permutation of \(xs\).
\end{lemma}

\begin{lstlisting}
Lemma qsort_perm xs : Permutation (qsort xs) xs.
\end{lstlisting}

\begin{lemma}[Sorted]
\(\mathsf{sorted}(\mathsf{qsort}(xs),\, 0,\, \lvert xs \rvert)\).
\end{lemma}

\begin{lstlisting}
Lemma qsort_sorted xs : slice_sorted (qsort xs) 0%nat (length xs).
\end{lstlisting}

One concrete run is computed in Coq.

\begin{lstlisting}
Lemma qsort_example : qsort [3%Z; 1%Z; 2%Z] = [1%Z; 2%Z; 3%Z].
\end{lstlisting}

\subsection{The partition invariant}

The scan is the one loop.
Its invariant relates the list at the start of the slice, \(xs_0\), to the list \(xs\) at an intermediate cursor \((i, j)\).
The pivot value sits at \(\mathit{bound}\) throughout the scan.
The final exchange of the pivot happens after the loop, so the invariant still has the pivot at the end of the slice.

\begin{lstlisting}
Record lomuto_inv (xs0 xs : list Z) (lo i j bound : nat) (pivot : Z) : Prop := {
  linv_len : length xs = length xs0;
  linv_perm : Permutation xs xs0;
  linv_lo : (lo <= i)%nat;
  linv_ij : (i <= j)%nat;
  linv_jb : (j <= bound)%nat;
  linv_lenb : (bound < length xs)%nat;
  linv_piv : xs !! bound = Some pivot;
  linv_out : forall k, (k < lo \/ bound < k)%nat -> xs !! k = xs0 !! k;
  linv_left : forall k x, (lo <= k < i)%nat -> xs !! k = Some x -> (x <= pivot)%Z;
  linv_mid : forall k x, (i <= k < j)%nat -> xs !! k = Some x -> (pivot < x)%Z
}.
\end{lstlisting}

Read as a picture of \([lo, \mathit{bound}]\):

\begin{itemize}
\item The length is unchanged, and the list is a permutation of the original.
\item The cursors satisfy \(lo \le i \le j \le \mathit{bound}\), and \(\mathit{bound}\) is in range.
\item Index \(\mathit{bound}\) still holds the pivot.
\item Every index outside \([lo, \mathit{bound}]\) is unchanged.
\item Every element of \([lo, i)\) is \(\le\) the pivot.
\item Every element of \([i, j)\) is \(>\) the pivot.
\item The segment \([j, \mathit{bound})\) has not been classified yet.
\end{itemize}

At the start, \(i = j = lo\), so both classified segments are empty.

\begin{lstlisting}
Lemma lomuto_inv_start xs lo bound pivot :
  (lo <= bound)%nat ->
  (bound < length xs)%nat ->
  xs !! bound = Some pivot ->
  lomuto_inv xs xs lo lo lo bound pivot.
\end{lstlisting}

One iteration preserves the invariant.
If \(j < \mathit{bound}\) and \(xs_j = v\), the next state is the swapped list with both cursors advanced when \(v \le \mathit{pivot}\), and the same list with only \(j\) advanced when \(v > \mathit{pivot}\).
In the first branch the old element at \(i\), which was in the unclassified or the greater segment, moves to \(j\), and \(v\) becomes the new end of the lesser prefix.
In the second branch \(v\) becomes the new end of the greater segment.

\begin{lstlisting}
Lemma lomuto_step xs0 xs lo i j bound pivot v :
  lomuto_inv xs0 xs lo i j bound pivot ->
  (j < bound)%nat ->
  xs !! j = Some v ->
  lomuto_inv xs0
    (if Z.leb v pivot then swap_list xs i j else xs)
    lo (if Z.leb v pivot then S i else i) (S j) bound pivot.
\end{lstlisting}

Running the fuelled scan from a state that satisfies the invariant, with fuel \(\mathit{bound} - j\), finishes in a state whose cursors are \((p, \mathit{bound})\).

\begin{lstlisting}
Lemma lomuto_go n xs0 xs lo i j bound pivot ys p :
  lomuto_inv xs0 xs lo i j bound pivot ->
  (n = bound - j)%nat ->
  lomuto n xs pivot i j bound = (ys, p) ->
  lomuto_inv xs0 ys lo p bound bound pivot.
\end{lstlisting}

The final exchange moves the pivot from \(\mathit{bound}\) to \(p\).
What partition establishes, for a non-empty in-range slice, is the predicate \texttt{part\_ok}.

\begin{lstlisting}
Definition part_ok (xs ys : list Z) (lo hi p : nat) (pivot : Z) : Prop :=
  length ys = length xs /\ Permutation ys xs /\ (lo <= p < hi)%nat /\
  ys !! p = Some pivot /\
  (forall k, (k < lo \/ hi <= k)%nat -> ys !! k = xs !! k) /\
  (forall k x, (lo <= k < p)%nat -> ys !! k = Some x -> (x <= pivot)%Z) /\
  (forall k x, (p < k < hi)%nat -> ys !! k = Some x -> (pivot < x)%Z).

Lemma partition_list_ok xs lo hi ys p :
  (lo < hi)%nat -> (hi <= length xs)%nat ->
  partition_list xs lo hi = (ys, p) ->
  exists pivot, part_ok xs ys lo hi p pivot.
\end{lstlisting}

So \(ys\) has the same length as \(xs\) and is a permutation of it, the returned index \(p\) lies in \([lo, hi)\), index \(p\) holds the pivot, indexes outside the slice are untouched, \([lo, p)\) is \(\le\) the pivot, and \((p, hi)\) is \(>\) the pivot.

The recursive calls sort those two sides.
A permutation of a side can rearrange it, so the comparison with the pivot has to be re-established after the recursive call.
The argument counts elements that fail the comparison.
\(\mathsf{count}(P, xs)\) is the number of elements of \(xs\) for which \(P\) returns true.

\begin{lstlisting}
Fixpoint count_if (P : Z -> bool) (xs : list Z) : nat :=
  match xs with
  | [] => 0
  | x :: xs => (if P x then 1%nat else 0%nat) + count_if P xs
  end.
\end{lstlisting}

If a slice of \(xs\) lies entirely in \(P\), a list \(ys\) of the same length that permutes \(xs\) and agrees with \(xs\) outside that slice still lies entirely in \(P\) on that slice.
The count of failures on the slice is zero on both sides, because the count is preserved by permutation and the outside regions contribute the same count.

\begin{lstlisting}
Lemma region_preserved (P : Z -> bool) (xs ys : list Z) (lo hi : nat) :
  length ys = length xs ->
  (lo <= hi <= length xs)%nat ->
  (forall k, (k < lo \/ hi <= k)%nat -> ys !! k = xs !! k) ->
  Permutation ys xs ->
  (forall k x, (lo <= k < hi)%nat -> xs !! k = Some x -> P x = true) ->
  forall k x, (lo <= k < hi)%nat -> ys !! k = Some x -> P x = true.
\end{lstlisting}

Applied to \(P(z) = (z \le \mathit{pivot})\) on the left of \(p\), and to \(P(z) = (\mathit{pivot} < z)\) on the right of \(p\), this is what lets the two sorted sides and the pivot in the middle combine into one sorted slice.
A slice of length at most one is sorted outright.

\begin{lstlisting}
Lemma slice_sorted_short xs lo hi :
  (hi <= lo + 1)%nat -> slice_sorted xs lo hi.
\end{lstlisting}

Two adjacent sorted slices, with every element of \([lo, p)\) at most the pivot and every element of \((p, hi)\) greater than the pivot, form one sorted slice.

\begin{lstlisting}
Lemma slice_sorted_join xs lo p hi pivot :
  (lo <= p)%nat -> (p < hi)%nat ->
  xs !! p = Some pivot ->
  (forall k x, (lo <= k < p)%nat -> xs !! k = Some x -> (x <= pivot)%Z) ->
  (forall k x, (p < k < hi)%nat -> xs !! k = Some x -> (pivot < x)%Z) ->
  slice_sorted xs lo p ->
  slice_sorted xs (S p) hi ->
  slice_sorted xs lo hi.
\end{lstlisting}

\subsection{The length, get, and set interface}

The core language has no interfaces and no generics, so the shared program is a Coq section.
Its parameters are a ghost type \(G\), a representation \(R\), and three expression bodies.

\begin{lstlisting}
Variable G : Type.
Variable R : G -> list Z -> val -> iProp.
Variable len_body get_body set_body : expr.
\end{lstlisting}

\(R(g, xs, a)\) says that the value \(a\) currently represents the list \(xs\), with ghost information \(g\).
For the array, \(G\) is unit and \(R\) ignores it.
For the array list, \(g\) is the pair of the capacity and the version.

Each body must be closed apart from its own parameter.
Length binds \texttt{"a"}.
Reading and writing bind \texttt{"args"}.
Substitution of any other name, whether a value or a stack location, leaves the body unchanged.
The six hypotheses have this shape.

\begin{lstlisting}
Hypothesis len_sub : forall x v, x <> "a" -> subst x v len_body = len_body.
Hypothesis len_subvar :
  forall x (l : loc), x <> "a" -> subst_var x l len_body = len_body.
Hypothesis get_sub : forall x v, x <> "args" -> subst x v get_body = get_body.
Hypothesis get_subvar :
  forall x (l : loc), x <> "args" -> subst_var x l get_body = get_body.
Hypothesis set_sub : forall x v, x <> "args" -> subst x v set_body = set_body.
Hypothesis set_subvar :
  forall x (l : loc), x <> "args" -> subst_var x l set_body = set_body.
\end{lstlisting}

The three specifications are the only facts the algorithm applies about the library.
Each one is stated at the call form the program builds: a closure applied to one struct of arguments.
The ghost \(g\) is the same in the precondition and the postcondition.
A write changes the list and leaves \(g\) alone.

\begin{lstlisting}
Hypothesis len_spec : forall g xs (a : val),
  R g xs a -*
  WP App (Rec None (Some "a") len_body) (Val a)
    {{ v, bi_pure (v = LitV (LitInt (Z.of_nat (length xs)))) * R g xs a }}.

Hypothesis get_spec : forall g xs (a : val) (i : nat) x,
  xs !! i = Some x ->
  R g xs a -*
  WP App (Rec None (Some "args") get_body)
       (Val (StructV [a; LitV (LitInt (Z.of_nat i))]))
    {{ v, bi_pure (v = LitV (LitInt x)) * R g xs a }}.

Hypothesis set_spec : forall g xs (a : val) (i : nat) y x,
  xs !! i = Some y ->
  R g xs a -*
  WP App (Rec None (Some "args") set_body)
       (Val (StructV [a; LitV (LitInt (Z.of_nat i)); LitV (LitInt x)]))
    {{ v, bi_pure (v = LitV LitUnit) * R g (<[i := x]> xs) a }}.
\end{lstlisting}

In mathematics, with the call written as the operation it prints to:

\begin{definition}[Library interface]
\[
\begin{aligned}
R(g, xs, a) &\vdash \Wpre\,(\mathsf{length}(a))
  \{\,v.\; \pure{v = \lvert xs \rvert} \sep R(g, xs, a)\,\},\\
xs_i = x \land R(g, xs, a) &\vdash \Wpre\,(\mathsf{get}(a, i))
  \{\,v.\; \pure{v = x} \sep R(g, xs, a)\,\},\\
xs_i = y \land R(g, xs, a) &\vdash \Wpre\,(\mathsf{set}(a, i, x))
  \{\,v.\; \pure{v = ()} \sep R(g,\, xs[i \mapsto x],\, a)\,\}.
\end{aligned}
\]
\end{definition}

The value that represents the container is a literal object.
Substituting into it does nothing, which is the predicate \texttt{rigid}.
An instance opens the representation far enough to see \(a = \mathsf{lit}(o)\) and then rebuilds it.
A literal satisfies \texttt{rigid} by computation, which is \texttt{lit\_rigid}.

\begin{lstlisting}
Definition rigid (v : val) : Prop :=
  (forall x w, subst_val x w v = v) /\
  (forall x (l : loc), subst_var_val x l v = v).
\end{lstlisting}

\subsection{The program}

Application in the core language is unary, so a call of two or three arguments builds a struct, binds it, and applies the library closure to that struct.
The printer erases the struct and prints \texttt{get(a, i)} or \texttt{a.get(i)}, according to which body was supplied.
\texttt{call2} binds the result.
\texttt{call3} is the call itself.
\texttt{call3\_seq} sequences the call before a continuation.

\begin{lstlisting}
Definition call2 (body : expr) (a b : expr) (dst : string) (k : expr) : expr :=
  Let (Some "pack") (Struct [] [a; b])
    (Let (Some dst) (App (Rec None (Some "args") body) (Var "pack")) k).

Definition call3 (body : expr) (a b c : expr) : expr :=
  Let (Some "pack") (Struct [] [a; b; c])
    (App (Rec None (Some "args") body) (Var "pack")).

Definition call3_seq (body : expr) (a b c k : expr) : expr :=
  Let (Some "pack") (Struct [] [a; b; c])
    (Let None (App (Rec None (Some "args") body) (Var "pack")) k).
\end{lstlisting}

Swap reads index \(i\) into \texttt{t}, reads index \(j\) into \texttt{u}, writes \texttt{u} at \(i\), and writes \texttt{t} at \(j\).

\begin{lstlisting}
Definition swap_body : expr :=
  Let (Some "a") (StructLoad (Var "args") 0)
  (Let (Some "i") (StructLoad (Var "args") 1)
  (Let (Some "j") (StructLoad (Var "args") 2)
    (call2 get_body (Var "a") (Var "i") "t"
      (call2 get_body (Var "a") (Var "j") "u"
        (call3_seq set_body (Var "a") (Var "i") (Var "u")
          (call3 set_body (Var "a") (Var "j") (Var "t"))))))).
\end{lstlisting}

\begin{lemma}[Swap]
If \(a\) is rigid, \(xs_i = x_i\), and \(xs_j = x_j\), then
\[
R(g, xs, a) \vdash
  \Wpre\,(\mathsf{swap}(a, i, j))
  \{\,v.\; \pure{v = ()} \sep R(g,\, \mathsf{swap}(xs, i, j),\, a)\,\}.
\]
\end{lemma}

\begin{lstlisting}
Lemma swap_spec g xs (a : val) (i j : nat) xi xj :
  rigid a ->
  xs !! i = Some xi -> xs !! j = Some xj ->
  R g xs a -*
  WP call_swap a i j
    {{ v, bi_pure (v = LitV LitUnit) * R g (swap_list xs i j) a }}.
\end{lstlisting}

The scan is the \texttt{while} of partition.
Its condition is \((j + 1) \le (hi - 1)\), that is \(j < hi - 1\).
The body reads index \(j\).
When the value is \(\le\) the pivot it swaps indexes \(i\) and \(j\) and advances \(i\).
It then advances \(j\).
The two cursors are mutable \texttt{var} bindings, allocated by \texttt{VarBind} from the initial value \(lo\).

\begin{lstlisting}
Definition scan (a hi pivot : expr) : expr :=
  While (BinOp LeOp (plus1 (Var "j")) (minus1 hi))
    (call2 get_body a (Var "j") "v"
      (Let None
        (If (BinOp LeOp (Var "v") pivot)
          (Let None (call3 swap_body a (Var "i") (Var "j"))
            (Assign "i" (plus1 (Var "i"))))
          unitv)
        (Assign "j" (plus1 (Var "j"))))).
\end{lstlisting}

After the two \texttt{var} bindings are opened, \(i\) and \(j\) reside in stack cells \(l_i\) and \(l_j\).
The loop specification is an induction on the number of steps remaining, \((hi - 1) - j\).
The postcondition is the Gallina scan: the list is \(\mathsf{lomuto}(n, xs, \mathit{pivot}, i, j, hi - 1)\), the cell \(l_i\) holds the returned boundary, and the cell \(l_j\) holds \(hi - 1\).
The invariant maintained at each iteration is \(\mathsf{lomutoInv}\) of the section above, instantiated at the current cursors.
The base case is the failing loop test, \(j = hi - 1\).

\begin{lstlisting}
Lemma wp_scan n g xs (a : val) (hi : nat) (pivot : Z) (i j : nat) (li lj : loc) :
  rigid a ->
  (0 < hi)%nat ->
  ((hi - 1) - j = n)%nat ->
  (j <= hi - 1)%nat ->
  (i <= j)%nat ->
  (hi - 1 < length xs)%nat ->
  stack_mapsto (StackId li) (nval i) -*
  stack_mapsto (StackId lj) (nval j) -*
  R g xs a -*
  WP scan_open a (Z.of_nat hi) pivot li lj
  {{ _, exists ys p,
      bi_pure (lomuto n xs pivot i j (hi - 1) = (ys, p)) *
      stack_mapsto (StackId li) (nval p) *
      stack_mapsto (StackId lj) (nval (hi - 1)) * R g ys a }}.
\end{lstlisting}

Partition reads the pivot at \(hi - 1\), allocates \(i\) and \(j\) at \(lo\), runs the scan, swaps the pivot into the returned boundary, and returns that boundary.

\begin{lstlisting}
Definition partition_body : expr :=
  Let (Some "a") (StructLoad (Var "args") 0)
  (Let (Some "lo") (StructLoad (Var "args") 1)
  (Let (Some "hi") (StructLoad (Var "args") 2)
    (call2 get_body (Var "a") (minus1 (Var "hi")) "pivot"
      (VarBind "i" (Var "lo")
        (VarBind "j" (Var "lo")
          (Let None (scan (Var "a") (Var "hi") (Var "pivot"))
            (Let None (call3 swap_body (Var "a") (Var "i") (minus1 (Var "hi")))
              (Var "i")))))))).
\end{lstlisting}

\begin{lemma}[Partition]
If \(a\) is rigid, \(lo < hi \le \lvert xs \rvert\), and \(xs_{hi-1} = \mathit{pivot}\), then
\[
\begin{aligned}
&R(g, xs, a) \vdash \Wpre\,(\mathsf{partition}(a, lo, hi))\\
&\qquad \{\,v.\; \exists\, ys, p.\;
  \pure{\mathsf{partition}(xs, lo, hi) = (ys, p)} \sep
  \pure{v = p} \sep R(g, ys, a)\,\}.
\end{aligned}
\]
\end{lemma}

\begin{lstlisting}
Lemma partition_spec g xs (a : val) (lo hi : nat) (pivot : Z) :
  rigid a ->
  (lo < hi)%nat ->
  (hi <= length xs)%nat ->
  xs !! (hi - 1)%nat = Some pivot ->
  R g xs a -*
  WP call_partition a lo hi
  {{ v, exists ys p,
      bi_pure (partition_list xs lo hi = (ys, p)) *
      bi_pure (v = nval p) * R g ys a }}.
\end{lstlisting}

The recursive procedure is named \texttt{qsort}.
Its test is \((lo + 1) \le (hi - 1)\), that is \(hi \ge lo + 2\).
The failing branch is unit, so the printed program has no \texttt{else}.
The succeeding branch partitions and calls itself on \([lo, p)\) and on \([p+1, hi)\).

\begin{lstlisting}
Definition qsort_body : expr :=
  Let (Some "a") (StructLoad (Var "args") 0)
  (Let (Some "lo") (StructLoad (Var "args") 1)
  (Let (Some "hi") (StructLoad (Var "args") 2)
    (If (BinOp LeOp (plus1 (Var "lo")) (minus1 (Var "hi")))
      (Let (Some "p")
        (App (Rec None (Some "args") partition_body)
          (Struct [] [Var "a"; Var "lo"; Var "hi"]))
        (Let None
          (App (Var "qsort") (Struct [] [Var "a"; Var "lo"; Var "p"]))
          (App (Var "qsort")
            (Struct [] [Var "a"; plus1 (Var "p"); Var "hi"]))))
      unitv))).
\end{lstlisting}

The specification follows the fuel of \texttt{qsort\_slice}.
With fuel at least \(hi - lo\), the call replaces the represented list by the sorted slice and leaves the ghost \(g\) unchanged.
Outside the slice the functional specification already says the elements are untouched, so the representation of the whole list is the specification of the slice.

\begin{lstlisting}
Lemma wp_qslice n g xs (a : val) (lo hi : nat) :
  rigid a ->
  (hi <= length xs)%nat ->
  (hi - lo <= n)%nat ->
  R g xs a -*
  WP call_qsort a lo hi {{ _, R g (qsort_slice n xs lo hi) a }}.
\end{lstlisting}

\texttt{sort} reads the length and calls \texttt{qsort} on \([0, n)\).

\begin{lstlisting}
Definition sort_body : expr :=
  Let (Some "n") (App (Rec None (Some "a") len_body) (Var "a"))
    (Let (Some "pack") (Struct [] [Var "a"; zero; Var "n"])
      (App (Rec (Some "qsort") (Some "args") qsort_body) (Var "pack"))).
\end{lstlisting}

\begin{theorem}[Sort]
If \(a\) is rigid, then
\[
R(g, xs, a) \vdash
  \Wpre\,(\mathsf{sort}(a))
  \{\,\_.\; R(g,\, \mathsf{qsort}(xs),\, a)\,\}.
\]
\end{theorem}

\begin{lstlisting}
Lemma sort_spec g xs (a : val) :
  rigid a ->
  R g xs a -*
  WP call_sort a {{ _, R g (qsort xs) a }}.
\end{lstlisting}

The proof is the slice specification at fuel \(\lvert xs \rvert\), on the slice \([0, \lvert xs \rvert)\).
The length read produces that upper bound.
A permutation preserves the length, so an array list's size, which is the length of \(xs\), is the same after sorting.

\subsection{The two instances}

The array instance takes \(G\) to be unit and \(R(\_, ys, b)\) to be \(\mathsf{isArray}(ys, b)\), the predicate of Section~\ref{sec:arrays}.
The three bodies are \texttt{length\_body}, \texttt{get\_body}, and \texttt{set\_body}.
Their closedness hypotheses are computed from those concrete terms.
Their specifications are the library lemmas \texttt{length\_spec}, \texttt{get\_spec}, and \texttt{set\_spec}, each applied in range.
The resulting statement mentions only the array predicate and \texttt{qsort}.

\begin{lstlisting}
Lemma array_qsort_spec xs (a : val) :
  is_array xs a -*
  WP call_sort length_body get_body set_body a
    {{ _, is_array (qsort xs) a }}.
\end{lstlisting}

The array-list instance takes \(G\) to be \(\mathbb{N} \times \mathbb{Z}\) and \(R((c, v), ys, b)\) to be \(\mathsf{isArrayList}(ys, c, v, b)\).
The three bodies are \texttt{size\_body}, \texttt{al\_get\_body}, and \texttt{al\_set\_body}.
The write specification of the library already returns the same capacity and the same version, and the interface requires exactly that.
Size is \(\lvert xs \rvert\) before the call and \(\lvert \mathsf{qsort}(xs) \rvert\) after it, and those lengths are equal.

\begin{lstlisting}
Lemma arraylist_qsort_spec xs cap ver (a : val) :
  is_arraylist xs cap ver a -*
  WP call_sort size_body al_get_body al_set_body a
    {{ _, is_arraylist (qsort xs) cap ver a }}.
\end{lstlisting}

\subsection{The printed programs}

\texttt{get} is a free function on \texttt{Array} and a method on \texttt{ArrayList}, and each generated file is compiled by itself, so the same three procedures are printed twice.
The array file carries \texttt{make}, \texttt{length}, \texttt{get}, \texttt{set}, \texttt{swap}, \texttt{partition}, \texttt{qsort}, and \texttt{sort}.
The array-list file carries \texttt{initArrayList}, \texttt{size}, \texttt{capacity}, \texttt{version}, \texttt{get}, \texttt{set}, \texttt{grow}, \texttt{add}, and the same three procedures.
A client fills the list with \texttt{add}, and \texttt{add} calls \texttt{grow}.
\texttt{sort} itself calls the length, the read, and the write.
On the array those calls are \texttt{length(a)}, \texttt{get(a, i)}, and \texttt{set(a, i, x)}.
On the array list they are \texttt{a.size()}, \texttt{a.get(i)}, and \texttt{a.set(i, x)}.

\lstinputlisting{generated/cj/qsort_array.cj}

\lstinputlisting{generated/cj/qsort_array_list.cj}

\section{Hash maps}
\label{sec:hashmap}

\texttt{theories/examples/hash\_map.v} lowers the monomorphic \texttt{Int64} core of Cangjie's \texttt{HashMap}.
Three \texttt{Array} fields hold tags, keys, and values; tag $0$ is empty, $1$ occupied, $2$ a tombstone.
The live size is stored in \texttt{mySize}.
The mathematical content is a finite map with no duplicate keys; \texttt{wf\_table} and the probe invariants tie the three arrays to that map.

The bucket index is \texttt{residue k cap}: add \texttt{cap} while the key is negative, then subtract \texttt{cap} while the representative is at least \texttt{cap}.
That is the remainder the source would obtain with a power-of-two mask.
Growth doubles the capacity when \texttt{size + size} reaches the capacity; the extracted program allocates fresh arrays (rehashing is still being wired to the \texttt{grow\_table} Gallina step).

\texttt{get} returns an \texttt{Int64Option}: a boolean \texttt{ok} and an integer \texttt{value}.
\texttt{contains} is the \texttt{ok} flag alone.
\texttt{remove} of a missing key returns zero and leaves the table unchanged.

\lstinputlisting{generated/cj/hash_map.cj}

\section{What is deliberately missing}

We deliberately omit concurrency from the semantics, so there is no invariant-opening rule and no fork.
Float literals exist, yet no float operator does.
Allocation of a positive word count is the rule used by the array library.
Freeing in the named logic rules remains the one-word rule.
Objects are not collected.
Pointer arithmetic is word-addressed.
Null, a freed address, and an address with no points-to fact are stuck rather than given a defined error value.

The type grammar is not enforced.
The adequacy theorem is the safety theorem for those programs whose specifications have been proved.

\section{Reading list}

The entries below are the sources this note relies on.
The formal bibliography follows the same order.

\paragraph{The language.}
The Cangjie 1.0 documentation~\cite{cangjie-docs} is the surface syntax from which this fragment is cut: \texttt{func}, \texttt{let} and \texttt{var}, \texttt{if} and \texttt{while}, \texttt{struct}, non-\texttt{open} \texttt{class}, and \texttt{CPointer}.
The performance-parity cut keeps unboxed structs, word-addressed raw blocks, and final classes, with generics, inheritance, and concurrency left as sugar that must lower into that cut.

\paragraph{Separation logic.}
Reynolds~\cite{reynolds-lics02} is the logic of disjoint ownership of mutable structure, which is the reading of \(\ptss\), \(\ptsr\), \(\ptsb\), and \(\ptso\) in this development.
O'Hearn, Reynolds, and Yang~\cite{ohearn-csl01} is the local-reasoning paper: a specification mentions only the cells the command uses, and the frame rule supplies the rest.
That is the justification of the framing steps in the examples and in \texttt{cons\_spec}.

\paragraph{Iris.}
The POPL 2015 paper~\cite{iris-popl15} introduces the resource algebra that the four ghost maps instantiate.
The ESOP 2017 paper~\cite{iris-essence} is the one to read for the program logic used here: weakest preconditions, evaluation contexts, the bind rule, and the lifting of a deterministic head step.
\emph{Iris from the ground up}~\cite{iris-ground-up} is the self-contained account of the same logic, including adequacy.
The Iris Proof Mode paper~\cite{iris-ipm} explains the \texttt{iApply}, \texttt{iDestruct}, and \texttt{iFrame} style of the Coq proofs.
The lecture notes~\cite{iris-notes} are the shortest path through all of that with exercises.
HeapLang, in the Iris 4.4.0 Coq development~\cite{iris-coq}, is the language whose \texttt{ectxi\_language} instance this one copies; \texttt{ghost\_map} is the library behind the points-to facts.

\paragraph{The proof assistant.}
The mechanization is checked with Coq 8.20.1~\cite{coq} and std++ 1.12.0~\cite{stdpp}.
std++ supplies the finite maps, the list lemmas, and the \texttt{decide} tactic used for the pure arithmetic steps.

\begin{thebibliography}{99}

\bibitem{cangjie-docs}
Cangjie documentation, version 1.0.0.
\url{https://docs.cangjie-lang.cn/en/docs/1.0.0/}.

\bibitem{reynolds-lics02}
John C. Reynolds.
Separation logic: A logic for shared mutable data structures.
In \emph{Proceedings of LICS}, pages 55--74. IEEE, 2002.

\bibitem{ohearn-csl01}
Peter W. O'Hearn, John C. Reynolds, and Hongseok Yang.
Local reasoning about programs that alter data structures.
In \emph{Proceedings of CSL}, volume 2142 of \emph{LNCS}, pages 1--19. Springer, 2001.

\bibitem{iris-popl15}
Ralf Jung, David Swasey, Filip Sieczkowski, Kasper Svendsen, Aaron Turon, Lars Birkedal, and Derek Dreyer.
Iris: Monoids and invariants as an orthogonal basis for concurrent reasoning.
In \emph{Proceedings of POPL}, pages 637--650. ACM, 2015.

\bibitem{iris-essence}
Robbert Krebbers, Ralf Jung, Ale\v{s} Bizjak, Jacques-Henri Jourdan, Derek Dreyer, and Lars Birkedal.
The essence of higher-order concurrent separation logic.
In \emph{Proceedings of ESOP}, volume 10201 of \emph{LNCS}, pages 696--723. Springer, 2017.

\bibitem{iris-ground-up}
Ralf Jung, Robbert Krebbers, Jacques-Henri Jourdan, Ale\v{s} Bizjak, Lars Birkedal, and Derek Dreyer.
Iris from the ground up: A modular foundation for higher-order concurrent separation logic.
\emph{Journal of Functional Programming}, 28:e20, 2018.

\bibitem{iris-ipm}
Robbert Krebbers, Amin Timany, and Lars Birkedal.
Interactive proofs in higher-order concurrent separation logic.
In \emph{Proceedings of POPL}, pages 205--217. ACM, 2017.

\bibitem{iris-notes}
Lars Birkedal and Ale\v{s} Bizjak.
Lecture notes on Iris: Higher-order concurrent separation logic.
\url{https://iris-project.org/pdfs/iris-lecture-notes.pdf}.

\bibitem{iris-coq}
The Iris Project.
The Iris Coq development, version 4.4.0, including HeapLang.
\url{https://gitlab.mpi-sws.org/iris/iris}.

\bibitem{coq}
The Coq Development Team.
The Coq proof assistant, version 8.20.1.
\url{https://coq.inria.fr/}.

\bibitem{stdpp}
The Iris Project.
coq-std++, version 1.12.0.
\url{https://gitlab.mpi-sws.org/iris/stdpp}.

\end{thebibliography}

\end{document}
